\documentclass[UTF8,a4paper,11pt,fontset=none]{ctexart}
\usepackage[top=22mm,bottom=22mm,left=22mm,right=22mm,headheight=16pt,headsep=7mm,footskip=11mm]{geometry}
\usepackage{fontspec}
\setmainfont{DejaVu Serif}
\setsansfont{DejaVu Sans}
\setmonofont[Scale=0.88]{DejaVu Sans Mono}
\setCJKmainfont[BoldFont=Noto Serif CJK SC Bold]{Noto Serif CJK SC}
\setCJKsansfont[BoldFont=Noto Sans CJK SC Bold]{Noto Sans CJK SC}
\setCJKmonofont{Noto Sans Mono CJK SC}
\usepackage{amsmath,amssymb,booktabs,longtable,array,calc}
\usepackage[table]{xcolor}
\definecolor{Ink}{HTML}{193347}
\definecolor{Accent}{HTML}{006D77}
\definecolor{Pale}{HTML}{EEF4F7}
\usepackage{fvextra}
\RecustomVerbatimEnvironment{verbatim}{Verbatim}{breaklines=true,breakanywhere=true,fontsize=\small,frame=single,rulecolor=\color{Accent},framesep=6pt}
\usepackage{xurl}
\usepackage{hyperref}
\hypersetup{unicode=true,colorlinks=true,linkcolor=Ink,urlcolor=Accent,pdftitle={LLM自动算法设计：学习与研究手册},pdfauthor={研究型学习手册},pdfsubject={指南核查、公开视频、源码、复现实验和24周学习计划}}
\urlstyle{same}
\usepackage{bookmark,fancyhdr,etoolbox,needspace}
\pagestyle{fancy}
\fancyhf{}
\fancyhead[L]{\sffamily\small\color{Ink}LLM 自动算法设计}
\fancyhead[R]{\sffamily\small\color{Ink}学习与研究手册}
\fancyfoot[L]{\sffamily\footnotesize 核验截止：2026-09-19}
\fancyfoot[R]{\sffamily\footnotesize\thepage}
\renewcommand{\headrulewidth}{0.35pt}
\renewcommand{\footrulewidth}{0pt}
\setlength{\parindent}{0pt}
\setlength{\parskip}{5.5pt}
\linespread{1.18}
\setlength{\emergencystretch}{3em}
\setcounter{secnumdepth}{0}
\setcounter{tocdepth}{1}
\ctexset{section={format=\Large\sffamily\bfseries\color{Ink},beforeskip=4pt,afterskip=15pt},subsection={format=\large\sffamily\bfseries\color{Ink},beforeskip=14pt,afterskip=7pt}}
\providecommand{\tightlist}{\setlength{\itemsep}{0pt}\setlength{\parskip}{0pt}}
\AtBeginEnvironment{longtable}{\small\renewcommand{\arraystretch}{1.15}\setlength{\LTpre}{8pt}\setlength{\LTpost}{10pt}}
\widowpenalty=10000
\clubpenalty=10000
\begin{document}
\begin{titlepage}
\thispagestyle{empty}
{\sffamily\color{Accent}\large 学习与研究手册\hfill 2026.09}
\vspace{26mm}

{\sffamily\bfseries\color{Ink}\fontsize{32}{43}\selectfont LLM 自动算法设计\par}
\vspace{5mm}
{\sffamily\color{Ink}\fontsize{22}{31}\selectfont 从指南纠错到独立复现\par}
\vspace{9mm}
{\color{Accent}\rule{\textwidth}{1.3pt}}
\vspace{7mm}
{\large 概念推导 · 视频与源码 · 实验规范 · 学习安排\par}
\vspace{15mm}
\begin{tabular}{@{}p{0.30\textwidth}p{0.30\textwidth}p{0.30\textwidth}@{}}
{\sffamily\bfseries\Huge 24周} & {\sffamily\bfseries\Huge 82条} & {\sffamily\bfseries\Huge 29项}\\[4mm]
主线学习计划 & 来源与资源记录 & 通过的教学测试\\
\end{tabular}\par
\vspace{16mm}
本手册围绕老师提供的入门指南，先修正概念、文献与代码问题，再以装箱、TSP和CVRP连接经典算法、LLM程序搜索及研究实验。目标是能够独立解释方法、运行小实验、识别证据边界，并设计可复核的研究。
\vfill
{\sffamily\small 审读对象：《LLM自动算法设计入门指南(1).docx》\par}
{\sffamily\small 检索核验截止：2026年9月19日\par}
{\sffamily\small 版本：1.0\quad Word与PDF共享正文；PDF由XeLaTeX编译\par}
\end{titlepage}
\pagenumbering{roman}
\tableofcontents
\clearpage
\pagenumbering{arabic}

\section{1
阅读方法与交付边界}\label{ux9605ux8bfbux65b9ux6cd5ux4e0eux4ea4ux4ed8ux8fb9ux754c}

这份手册以老师提供的《LLM自动算法设计入门指南(1).docx》为审读对象，目标不是把里面的名词都背下来，而是让你能够独立回答四个问题：\textbf{要设计什么算法；模型究竟改变了什么；改进由谁验证；证据能支持多大范围的结论。}
本册同时提供学习顺序、公开视频与源码入口、实验设计、进度安排和可以直接运行的纠错练习。

原指南的研究方向值得认真学习，但它不是一份可以逐条照抄执行的实验协议。正文中存在文献标识符错配、框架身份混淆、过时的模型推荐和不可直接运行的示例。这里分别处理``方向是否有价值''``资料是否准确''``代码是否可用''，不让任何一项代替另外两项。

\textbf{证据口径。}
检索核验截至2026年9月19日。原指南依据该文件可检索的正文逐节审读；原始Word字节导出未成功，因此不声称审查了嵌入对象、批注、修订历史或原始排版。外部资源以论文作者、出版机构、会议、大学课程和项目官方仓库为主；社区内容用于定位资源和解决工程问题，不替代一手证据。这是一份有检索记录的研究型学习设计，不是宣称穷尽全部平台的系统综述。

\textbf{执行口径。}
本次实际执行了附带的遗传算法、微型TSP、在线装箱及验证器练习，29项测试通过。没有在本次会话中完成EoH、FunSearch、ReEvo或大模型微调的整套论文复现；它们在本册中分别标为``已核验资料''``计划实验''。公开视频核验的是来源、主题与入口，不声称全部完整观看，也不承诺每个校园网和手机客户端均可播放。

\textbf{时间口径。}
主线按24周、每周14小时设计，共336小时；不足的编程或数学基础另加0---8周，论文研究另加8---12周。这里的小时是建议投入，不是课程视频时长，也不是保证掌握知识的工期。先通过能力检查，再决定是否压缩或补课。

\textbf{论文口径。}
``顶刊水准''在这里落实为可追溯引用、可检验假设、强且公平的基线、独立评估、预算核算、统计不确定性、明确的失败边界。漂亮排版、模型数量、代码体量和结果最好的一次，都不能替代这些要求。\href{https://jmlr.org/papers/v22/20-303.html}{{[}P22{]}}

\subsection{1.1
先走哪一条阅读路线}\label{ux5148ux8d70ux54eaux4e00ux6761ux9605ux8bfbux8defux7ebf}

第一次通读先看第2---4章，把指南纠错和核心问题看清，再看第10章实验协议与第11章时间安排。准备动手时读第7---9章；每完成一个练习再回看论文精读卡。第12章的研究候选不应在第一周就变成已经确定的``创新点''。所有资源编号均可点击；完整入口、核验状态与备注另见附录和随包的CSV/JSON目录。

\section{2
原指南逐节审读：保留主线，重排依赖}\label{ux539fux6307ux5357ux9010ux8282ux5ba1ux8bfbux4fddux7559ux4e3bux7ebfux91cdux6392ux4f9dux8d56}

\subsection{2.1
原指南的八个部分，分别需要学到什么程度}\label{ux539fux6307ux5357ux7684ux516bux4e2aux90e8ux5206ux5206ux522bux9700ux8981ux5b66ux5230ux4ec0ux4e48ux7a0bux5ea6}

\begin{longtable}[]{@{}
  >{\raggedright\arraybackslash}p{(\columnwidth - 4\tabcolsep) * \real{0.23}}
  >{\raggedright\arraybackslash}p{(\columnwidth - 4\tabcolsep) * \real{0.35}}
  >{\raggedright\arraybackslash}p{(\columnwidth - 4\tabcolsep) * \real{0.42}}@{}}
\toprule\noalign{}
\begin{minipage}[b]{\linewidth}\raggedright
原指南部分
\end{minipage} & \begin{minipage}[b]{\linewidth}\raggedright
应保留的学习目标
\end{minipage} & \begin{minipage}[b]{\linewidth}\raggedright
本册的调整与验收
\end{minipage} \\
\midrule\noalign{}
\endhead
\bottomrule\noalign{}
\endlastfoot
1 引言 & 理解自动算法设计，而非只会调用聊天模型 &
能区分生成一个解、生成一个算法、调参数、选算法 \\
2 基础知识 & EA、NAS、程序合成、数学、Python、LLM &
分成必须会用、必须理解、方向确定后深入三层 \\
3 学习步骤 & 从基础到复现再到独立研究 &
用能力关卡替代``刷满200题''和每日7---8小时 \\
4 代码与GitHub & 阅读真实项目并接通评估闭环 &
核对项目身份、真实API、版本与可执行范围 \\
5 关键论文 & 建立概念谱系、理解方法与证据 &
修正错引；把经典、平台、模型论文与近期预印本分开 \\
6 实践项目 & 从小实验走到可复现实验 &
保留六类目标，不同时启动六个大工程 \\
7 学习路线图 & 形成连续的产出与复盘 &
每周一个可检查成果，失败则补前置条件 \\
8 链接汇总 & 能持续找到可信资源 &
建资源登记表，不以链接数或收藏量表示掌握程度 \\
\end{longtable}

原指南没有把应用限定为农业、CVRP或覆盖规划。本册后面的农业任务规划是结合你的研究语境给出的迁移建议，不冒充原指南的硬性规定。自动算法设计应是学习主线；TSP与装箱先充当容易检查的教学任务，最终研究任务等完成基础复现后再选。

\subsection{2.2
三种学习深度，不把所有前置课都当成门槛}\label{ux4e09ux79cdux5b66ux4e60ux6df1ux5ea6ux4e0dux628aux6240ux6709ux524dux7f6eux8bfeux90fdux5f53ux6210ux95e8ux69db}

\textbf{必须亲手实现。}
输入检查、随机种子与实验记录、至少一个图搜索算法、TSP路线代价与可行性验证、最近邻与2-opt、First
Fit/Best
Fit、一个最小遗传算法、数据分组划分、简单统计汇总。这里的要求是离开原教程仍能实现，并能设计反例。

\textbf{必须解释并做小实验。}
Attention的矩阵形状与因果掩码，预训练/监督微调/强化学习的区别，程序合成的规格与验证，GA/PSO/DE的更新对象，NAS的搜索空间与评估代价，RL的状态/动作/奖励，模仿学习的分布偏移。掌握这些，不等于需要先训练一个大型语言模型。

\textbf{确定方向后再深入。}
大模型权重微调、复杂NAS、多目标选择器、完整求解器自动生成、跨任务迁移、对抗式实例生成、多机大规模实验。它们不是``越多越好''的模块，而是针对具体瓶颈的手段。

\subsection{2.3
原指南六类项目如何处理}\label{ux539fux6307ux5357ux516dux7c7bux9879ux76eeux5982ux4f55ux5904ux7406}

GA/PSO项目保留为评分和预算意识训练；代码模型生成排序函数保留为接口与正确性练习，但不绑定某个历史模型；EoH复现是主线项目，不把微调列为先决条件；NAS保留一个小型完整实验作为概念分支；综合系统只围绕一个经过验证的评估器搭建；研究项目最后才讨论创新。这既覆盖指南，也避免一年都在补课或安装框架。

\section{3
必须先修正的事实、引文与代码问题}\label{ux5fc5ux987bux5148ux4feeux6b63ux7684ux4e8bux5b9eux5f15ux6587ux4e0eux4ee3ux7801ux95eeux9898}

\subsection{3.1
分类与论文身份}\label{ux5206ux7c7bux4e0eux8bbaux6587ux8eabux4efd}

\textbf{四角色分类被改写了。}
指南把LLM角色写成Optimizer、Programmer、Evolver、Designer，并归于Liu等综述；本次核验的综述v5使用的是optimizers、predictors、extractors、designers。你可以为了教学自定义``程序员''``进化算子''等角色，但必须声明这是自己的功能分类，不能当作原文术语。进一步说，``EoH用了进化''不等于模型必然属于某一个互斥的单一类别。\href{https://arxiv.org/abs/2410.14716v5}{{[}P01{]}}

\textbf{FunSearch、AlphaTensor和AlphaEvolve是三件不同的事。}
FunSearch用语言模型提出程序、通过评估和程序搜索改善数学构造及在线装箱启发式；AlphaTensor以强化学习搜索矩阵乘法算法；AlphaEvolve是更一般的进化式代码改进系统。不能把AlphaTensor的标题、作者、年份和矩阵乘法贡献贴到FunSearch名下。\href{https://deepmind.google/blog/funsearch-making-new-discoveries-in-mathematical-sciences-using-large-language-models/}{{[}P03{]}}
\href{https://arxiv.org/abs/2506.13131}{{[}P06{]}}
\href{https://deepmind.google/blog/discovering-novel-algorithms-with-alphatensor/}{{[}P25{]}}

\textbf{综述、算法论文、平台论文、模型论文不能相互替代。}
EoH论文回答的是怎样搜索启发式；LLM4AD回答的是怎样组织方法、任务和模型接口；CodeGen/StarCoder/CodeLlama回答的是代码模型的训练与能力；HumanEval主要测功能正确性。把它们统一排成``最新SOTA十篇''，会掩盖它们测量的对象完全不同。\href{https://arxiv.org/abs/2401.02051}{{[}P02{]}}
\href{https://arxiv.org/abs/2412.17287v2}{{[}P05{]}}
\href{https://arxiv.org/abs/2203.13474}{{[}P13{]}}
\href{https://arxiv.org/abs/2308.12950}{{[}P19{]}}
\href{https://arxiv.org/abs/2305.06161}{{[}P20{]}}
\href{https://github.com/openai/human-eval}{{[}C18{]}}

\subsection{3.2
文献标识符核验表}\label{ux6587ux732eux6807ux8bc6ux7b26ux6838ux9a8cux8868}

下表核对的是指南给出的``题名---链接''对应关系。链接能打开不代表引文正确；找不到某个条目也不等于证明世界上没有同名研究。

\begin{longtable}[]{@{}
  >{\raggedright\arraybackslash}p{(\columnwidth - 4\tabcolsep) * \real{0.23}}
  >{\raggedright\arraybackslash}p{(\columnwidth - 4\tabcolsep) * \real{0.35}}
  >{\raggedright\arraybackslash}p{(\columnwidth - 4\tabcolsep) * \real{0.42}}@{}}
\toprule\noalign{}
\begin{minipage}[b]{\linewidth}\raggedright
指南条目或标识符
\end{minipage} & \begin{minipage}[b]{\linewidth}\raggedright
核验结果
\end{minipage} & \begin{minipage}[b]{\linewidth}\raggedright
阅读与引用处理
\end{minipage} \\
\midrule\noalign{}
\endhead
\bottomrule\noalign{}
\endlastfoot
自动算法设计综述 & 应使用A Systematic Survey on Large Language Models
for Algorithm Design，2410.14716 & 固定v5，按其实际分类阅读
\href{https://arxiv.org/abs/2410.14716v5}{{[}P01{]}} \\
代码LLM综述：2305.08857 & 实为COWPEA，一篇选举方法论文
\href{https://arxiv.org/abs/2305.08857}{{[}A01{]}} &
替换为Jiang等代码生成综述
\href{https://arxiv.org/abs/2406.00515}{{[}P16{]}} \\
程序合成综述：2308.12345 & 实为早期引力修正研究
\href{https://arxiv.org/abs/2308.12345}{{[}A02{]}} &
不保留错引；先学程序合成课程，再读Austin等研究论文
\href{https://people.csail.mit.edu/asolar/SynthesisCourse/index.htm}{{[}V18{]}}
\href{https://arxiv.org/abs/2108.07732}{{[}P17{]}} \\
NAS综述：1908.09791 & 实为Once-for-All，并非Elsken综述
\href{https://arxiv.org/abs/1908.09791}{{[}A03{]}} & 综述改为1808.05377
\href{https://arxiv.org/abs/1808.05377}{{[}P15{]}} \\
EA与DL综述：2301.02345 & 实为手征U(1)规范理论中的轴子研究
\href{https://arxiv.org/abs/2301.02345}{{[}A04{]}} &
原条目撤出必读；不强配一个未经核验的替代标题 \\
FunSearch：2205.00527 & 实为Schmidt问题的有限类比
\href{https://arxiv.org/abs/2205.00527}{{[}A05{]}}；所写标题又属于AlphaTensor
& 分别引用FunSearch与AlphaTensor
\href{https://deepmind.google/blog/funsearch-making-new-discoveries-in-mathematical-sciences-using-large-language-models/}{{[}P03{]}}
\href{https://deepmind.google/blog/discovering-novel-algorithms-with-alphatensor/}{{[}P25{]}} \\
AlphaEvolve：2501.12345 & 实为The doubly librating Plutinos
\href{https://arxiv.org/abs/2501.12345}{{[}A06{]}} & 改为2506.13131
\href{https://arxiv.org/abs/2506.13131}{{[}P06{]}} \\
LLM4AD：2502.23456 & 本次返回404
\href{https://arxiv.org/abs/2502.23456}{{[}A07{]}} & 改为2412.17287
\href{https://arxiv.org/abs/2412.17287v2}{{[}P05{]}} \\
DARTS：1806.09055 & 题名匹配 & 保留，按NAS分支阅读
\href{https://arxiv.org/abs/1806.09055}{{[}P14{]}} \\
AlphaCode：2203.07814 & 题名匹配 &
保留为代码合成背景，不当作优化启发式评估
\href{https://arxiv.org/abs/2203.07814}{{[}P12{]}} \\
CodeGen：2203.13474 & 题名匹配 & 保留为历史模型论文
\href{https://arxiv.org/abs/2203.13474}{{[}P13{]}} \\
EoH：2401.02051 & 题名匹配 & 主线精读；核对真正的五种提示策略
\href{https://arxiv.org/abs/2401.02051}{{[}P02{]}} \\
CodeLlama / StarCoder / WizardCoder & 2308.12950 / 2305.06161 /
2306.08568对应论文存在 & 保留历史脉络，不直接称2026最新推荐
\href{https://arxiv.org/abs/2308.12950}{{[}P19{]}}
\href{https://arxiv.org/abs/2305.06161}{{[}P20{]}}
\href{https://arxiv.org/abs/2306.08568}{{[}P21{]}} \\
``Self-Evolving Code Generation''等五个泛化标题 &
指南缺少足够的作者、会场、标识符信息 & 暂不计为五篇已识别的必读论文 \\
\end{longtable}

指南把WizardCoder直接标为``Microsoft''，不足以作为规范作者信息；按论文作者署名引用即可。类似``支持商业使用''``最新70B''这样的判断，也不能由旧教程一句话替代当前模型卡、许可证与发布说明。本册不依据这些表述要求你下载历史大模型。\href{https://arxiv.org/abs/2306.08568}{{[}P21{]}}

指南中的``DSP（Dataflow Synthesis
Planner）''没有给出可追溯的作者、论文或仓库；本次精确词组检索未定位到可确认的对应项目。因此暂列身份未确认，不能根据缩写擅自补成另一个DSP/DSPy项目。

\subsection{3.3
三个仓库，不要再装错}\label{ux4e09ux4e2aux4ed3ux5e93ux4e0dux8981ux518dux88c5ux9519}

LLM4AlgorithmDesign是综述配套文献集合；EoH是方法实现；LLM4AD是包含多种方法与任务的通用平台。指南给LLM4AlgorithmDesign写出的目录树和\path{examples/run_eoh.py}安装命令，不能当作已核实的仓库文件。平台身份应从官方README与实际源码确定。\href{https://github.com/FeiLiu36/LLM4AlgorithmDesign}{{[}C01{]}}
\href{https://github.com/FeiLiu36/EoH}{{[}C02{]}}
\href{https://github.com/Optima-CityU/LLM4AD}{{[}C03{]}}

更重要的是，\textbf{真实仓库也会换API。} 本次读取EoH
README时，主分支已使用\texttt{LLMConfig\ /\ EoH\ /\ BaseProblem}；旧的\texttt{Paras\ /\ eoh.EVOL}对应旧接口，README指向v0.1。LLM4AD则有自己的\texttt{HttpsApi}、\texttt{OBPEvaluation}和\texttt{llm4ad.method.eoh}接口。不要把两个项目的代码段拼在一起运行。当前主分支能用，也不等于它与2024论文设置相同。\href{https://github.com/FeiLiu36/EoH}{{[}C02{]}}
\href{https://github.com/Optima-CityU/LLM4AD}{{[}C03{]}}

\subsection{3.4
两段示例代码的实质问题}\label{ux4e24ux6bb5ux793aux4f8bux4ee3ux7801ux7684ux5b9eux8d28ux95eeux9898}

\textbf{遗传算法：先繁殖后评分的顺序错误。}
原示例给旧种群计算fitness，然后创建fitness为None的子代，用子代替换种群，随即对None取max并格式化。这里不是``效果可能不好''，而是执行顺序会触发错误。附带练习独立复现了这个操作顺序，实际得到TypeError；没有冒称执行了无法导出的完整原文件。

\textbf{方向与边界也需要明确。}
原示例用max选父代；若直接把Sphere函数作为目标，则它在最大化本应最小化的值。只有明确取负适应度或更换选择方向才正确。单点交叉在维度为1时没有合法切点；非有限返回值、空输入、种群规模和停止条件也没有完整处理。精英保留不是GA的必要定义，但若你承诺``历史最好值单调不变差''，就必须显式保留并测试这个性质。\href{https://deap.readthedocs.io/en/master/tutorials/basic/part2.html}{{[}C12{]}}

\textbf{代码模型封装：生成文本不等于得到可执行候选。}
原示例把提示词和新生成内容一起解码，使用总长度限制而不是明确的新token预算，未完整处理模型设备、依赖、截断、代码提取、接口校验和执行隔离。纠正方向是：固定输入输出协议、保存原始响应、只提取合法候选、检查签名、独立评估；不是把生成结果交给\texttt{exec}就算完成。FP16、设备映射和依赖是否合适取决于具体模型与环境，不能简单判定为所有环境都错。\href{https://huggingface.co/learn/llm-course/zh-CN/chapter1/1}{{[}V12{]}}
\href{https://github.com/openai/human-eval}{{[}C18{]}}

\subsection{3.5
EoH五种策略的正确理解}\label{eohux4e94ux79cdux7b56ux7565ux7684ux6b63ux786eux7406ux89e3}

原论文把进化提示分成探索E1/E2和修改M1/M2/M3：E1倾向产生与父代不同的启发式；E2提取父代共性并据此发展新启发式；M1修改一个启发式；M2调整其参数；M3简化冗余部分。初始化另计，不能把五种策略随意改写成``交叉、变异、反思、学习、初始化''。这些是启发式操作的语义，不是保证成功的定理。\href{https://arxiv.org/abs/2401.02051}{{[}P02{]}}

\section{4
建立完整的概念模型}\label{ux5efaux7acbux5b8cux6574ux7684ux6982ux5ff5ux6a21ux578b}

\subsection{4.1
实例、解、算法、算法设计器：四层对象}\label{ux5b9eux4f8bux89e3ux7b97ux6cd5ux7b97ux6cd5ux8bbeux8ba1ux5668ux56dbux5c42ux5bf9ux8c61}

设\(x\)是一份问题实例，\(y\)是一份候选解，\(h\)是一个求解算法，\(M\)是算法设计过程。普通求解计算\(y=h(x)\)；自动算法设计计算\(h=M(D,B)\)，其中\(D\)是可用于设计的实例集合，\(B\)是设计预算。它们的输入、输出和泛化要求都不同。

以TSP为例，四个坐标点是实例，``0→1→2→3→0''是解，``每次选最近未访问点''是算法，``生成并筛选多个选点规则''是算法设计器。让LLM为一张图输出一条路线，不等于它已经设计出可迁移到其他图的算法；让LLM照着说明实现最近邻，也不自动成为新算法。\href{https://arxiv.org/abs/2410.14716v5}{{[}P01{]}}
\href{https://arxiv.org/abs/2401.02051}{{[}P02{]}}

\begin{longtable}[]{@{}
  >{\raggedright\arraybackslash}p{(\columnwidth - 4\tabcolsep) * \real{0.23}}
  >{\raggedright\arraybackslash}p{(\columnwidth - 4\tabcolsep) * \real{0.35}}
  >{\raggedright\arraybackslash}p{(\columnwidth - 4\tabcolsep) * \real{0.42}}@{}}
\toprule\noalign{}
\begin{minipage}[b]{\linewidth}\raggedright
易混概念
\end{minipage} & \begin{minipage}[b]{\linewidth}\raggedright
改变的对象
\end{minipage} & \begin{minipage}[b]{\linewidth}\raggedright
必须回答的问题
\end{minipage} \\
\midrule\noalign{}
\endhead
\bottomrule\noalign{}
\endlastfoot
算法选择 & 从既有算法集合中选一个 & 特征获取和选择的代价是否计入？ \\
算法配置 & 固定算法的参数 & 参数搜索是否比随机搜索更有效？ \\
算法生成 & 规则、表达式、代码或结构 &
生成内容是不是能在新实例上反复使用？ \\
神经求解器 & 模型权重与策略 &
学到的是直接构造解，还是帮助另一个求解器？ \\
超启发式 & 选择或生成启发式的高层方法 &
决策层在哪，底层算子是否固定？ \\
程序合成 & 满足规格的程序 &
规格是自然语言、测试用例，还是形式化逻辑？ \\
\end{longtable}

这些范围可以重叠，但实验必须说清楚重叠在哪里。模型既提议参数又改写代码，就不能把其收益全部解释成``选算法更聪明''。\href{https://arxiv.org/abs/2410.14716v5}{{[}P01{]}}
\href{https://arxiv.org/abs/2108.07732}{{[}P17{]}}
\href{https://2023.automl.cc/program/tutorials/survey-of-methods/}{{[}V22{]}}

\subsection{4.2
双层优化：外层设计算法，内层执行算法}\label{ux53ccux5c42ux4f18ux5316ux5916ux5c42ux8bbeux8ba1ux7b97ux6cd5ux5185ux5c42ux6267ux884cux7b97ux6cd5}

对于最小化任务，可以把设计目标写成：

\[
h^\star\in\arg\min_{h\in\mathcal H}
\mathbb E_{x\sim\mathcal D}\bigl[\ell(h,x;b)\bigr],
\qquad \mathrm{Cost}_{\mathrm{design}}\le B.
\]

这里\(b\)是每个求解算法执行一个实例时的预算，\(B\)是搜索算法本身的预算。前者可以是固定CPU时间或函数评价次数，后者包括模型调用、候选评估、修复和选择开销。用五小时设计出一个启发式，再让它处理一万个实例，与每个实例都请求LLM重新求解，是不同的部署模式。

这个表达式并不意味着我们知道真实分布\(\mathcal D\)，也不保证找到全局最优。实践中只得到有限数据上的估计值。需要独立验证集选择候选，再用封存测试集估计泛化。代码越自由、搜索越久，越有机会利用评估器漏洞或适应某一批样本；``搜索更多''因此同时增加能力和过拟合风险。

\subsection{4.3
从评分函数开始，不必直接生成整套求解器}\label{ux4eceux8bc4ux5206ux51fdux6570ux5f00ux59cbux4e0dux5fc5ux76f4ux63a5ux751fux6210ux6574ux5957ux6c42ux89e3ux5668}

一个最小的在线装箱接口可以写成：\texttt{score(item,\ remaining\_capacities)\ -\textgreater{}\ scores}。固定外层程序负责剔除装不下的箱子、处理新箱、并列规则和更新状态；候选函数只负责给合法选择排序。这样，方法差异集中在一个可以测量的部件上。

随后再逐步扩展到``局部搜索的惩罚函数''``算子选择规则''``多个协作函数''``完整求解器''。表达能力越大，验证、调试和归因越困难。当前EoH仓库已出现函数、多函数和类等不同设计范围，说明单函数不是领域上限；但它仍然是最合适的第一份因果清晰的教学实验之一。\href{https://github.com/FeiLiu36/EoH}{{[}C02{]}}

\subsection{4.4
正确性、可行性、质量、安全性必须分开}\label{ux6b63ux786eux6027ux53efux884cux6027ux8d28ux91cfux5b89ux5168ux6027ux5fc5ux987bux5206ux5f00}

``代码可以导入''只排除了部分语法错误；``返回一个列表''不保证它是合法路线；``路线合法''不保证代价算对；``代价较低''不保证没有读测试答案；``在子进程运行''也不保证安全。验证至少分为接口检查、结果可行性、独立目标重算和运行隔离四层。\href{https://github.com/openai/human-eval}{{[}C18{]}}

用HumanEval等任务练习代码生成很有帮助，但优化实验还要报告路线代价、求解时间、失败率和分布外性能。pass@k测的是采样得到正确程序的机会，不能替代启发式的优化质量，更不能忽略k次生成的预算。\href{https://github.com/openai/human-eval}{{[}C18{]}}

\section{5
数学与算法：沿着三个可计算的例子理解}\label{ux6570ux5b66ux4e0eux7b97ux6cd5ux6cbfux7740ux4e09ux4e2aux53efux8ba1ux7b97ux7684ux4f8bux5b50ux7406ux89e3}

\subsection{5.1
在线装箱：最小的``算法设计''练习}\label{ux5728ux7ebfux88c5ux7bb1ux6700ux5c0fux7684ux7b97ux6cd5ux8bbeux8ba1ux7ec3ux4e60}

箱子容量为1，物品按0.6、0.6、0.4、0.4依次到达。Best
Fit选择放入后剩余容量最小的可行箱子；Next
Fit只尝试当前箱子。前者最终用2个箱子，后者用3个。附带代码计算并验证了这个差异。重要的不是记住谁赢，而是知道在线算法此刻不能读取未来物品；把未来物品喂给评分函数，就悄悄改变了问题。

令剩余容量为\(r_j\)、当前物品为\(a\)，Best
Fit可以在\(r_j\ge a\)的合法箱子集合里最小化\(r_j-a\)。若评估器采取最大分数选择，可令\(s_j=-(r_j-a)\)。最大化评分与最小化实际目标不矛盾，前提是符号、掩码与并列规则都写清楚。

若以\(LB=\lceil\sum_i a_i/C\rceil\)作为简单体积下界，使用箱数\(N_h\)可报告\((N_h-LB)/LB\)。必须称``相对该下界的差距''，不能称``最优差距''，因为这个下界未必紧。空实例需单独定义，不通过随意加一个很小的数掩盖分母为零。

第一轮实验只比较Next Fit、First Fit、Best
Fit与同接口候选；保留原始物品顺序、容量、分布和随机种子。后续再增加接近半箱的困难物品、容量尺度变化和未见分布。它的价值在于评估器便宜、可行性好检查，不代表它天然最适合农业研究。

\subsection{5.2
TSP：精确小实例与规模化启发式}\label{tspux7cbeux786eux5c0fux5b9eux4f8bux4e0eux89c4ux6a21ux5316ux542fux53d1ux5f0f}

对\(n\)个节点及代价矩阵\(D\)，一条合法巡回访问每个节点恰好一次并回到起点：

\[
L(\pi)=\sum_{i=0}^{n-1}D_{\pi_i,\pi_{(i+1)\bmod n}}.
\]

四个单位正方形顶点按周边访问，长度为4；选择一条含两条对角线的交叉路线，长度为\(2+2\sqrt2\approx4.828427\)。这里最优值可用枚举核对；一般大实例则不能把启发式找到的最好值自动当成已证明的最优值。

你应实现：固定起点的穷举小实例求解器、最近邻、2-opt和独立路线验证器。穷举仅用于\(n\le9\)的教学检查，不能拿它当大规模基线。对于对称TSP，2-opt可用边替换快速计算增量；对于非对称代价，片段反转会改变内部边方向，照抄对称公式可能算错。因此先全量重算验证，再优化增量更新。

比例缩放、节点重编号和欧氏平移旋转是有用的测试。欧氏坐标旋转不会改变距离矩阵，但按坐标轴硬编码的规则可能变；存在并列距离时，不应强迫算法返回完全相同的节点序列，而应检查代价、可行性和事先声明的并列策略。

\subsection{5.3
CVRP：从节点顺序走向多机器人任务分配}\label{cvrpux4eceux8282ux70b9ux987aux5e8fux8d70ux5411ux591aux673aux5668ux4ebaux4efbux52a1ux5206ux914d}

CVRP给每个客户需求\(q_i\)，每条车路线容量不超过\(Q\)，从仓库出发并返回。不同车辆之间既要分配任务，也要安排顺序。因此``把TSP多跑几遍''不是一般正确的CVRP建模方法。容量、车辆数量是否固定、是否允许拆分需求、是否有时间窗，都要在问题定义中写明。\href{https://github.com/PyVRP/PyVRP}{{[}C13{]}}
\href{https://github.com/vidalt/HGS-CVRP}{{[}C14{]}}
\href{https://developers.google.com/optimization/routing/vrp}{{[}C15{]}}

为教学目的，先做同质车辆、单仓库、确定性非负代价、不可拆分需求、无时间窗的CVRP。用成熟求解器建立强基线，再考虑自己的候选。引入农业属性时，不要立即叠加所有约束；先回答某个约束是否真的改变可行域、代价或最优决策，再决定是否加入。

一个可复核的微型例子：容量\(Q=10\)，三个客户需求分别为6、6、4，则总需求给出的车辆数下界为2，但哪些点同车仍由路线代价决定。若一个客户需求为11且禁止拆分，则实例本身不可行；这是数据/模型层问题，不能被算法返回空路线掩盖。

\subsection{5.4
EA、PSO、DE：看清``被更新的东西''}\label{eapsodeux770bux6e05ux88abux66f4ux65b0ux7684ux4e1cux897f}

GA可更新数值向量、排列、表达式或程序；交叉和变异必须尊重表示。普通实数交叉不能直接保证生成合法TSP排列。PSO通常以位置和速度更新候选数值向量；DE用若干个体的差分构造试验向量。它们不是因为名字不同就应全部塞进LLM框架。

入门时可用同一个Sphere目标\(f(z)=\sum_j z_j^2\)比较GA、随机搜索和一个DE实现，但必须匹配函数评价次数、初始化范围与维度。曲线横轴先用评价次数，再补墙钟时间；只比较``迭代了30代''会忽略不同种群和算法每代工作量不相同。\href{https://github.com/DEAP/deap}{{[}C11{]}}
\href{https://deap.readthedocs.io/en/master/tutorials/basic/part2.html}{{[}C12{]}}

在LLM算法设计里，个体由数值向量变成程序，适应度由一次函数值变成一批问题实例上的求解表现。这个变化比``把随机变异改成提示词''更根本：评估开始涉及程序错误、运行时间、泛化和安全。

\subsection{5.5
Attention与代码生成：只保留真正要用的数学}\label{attentionux4e0eux4ee3ux7801ux751fux6210ux53eaux4fddux7559ux771fux6b63ux8981ux7528ux7684ux6570ux5b66}

单头注意力可以写成：

\[
\operatorname{Attention}(Q,K,V)=
\operatorname{softmax}\!\left(\frac{QK^\top}{\sqrt{d_k}}+M\right)V.
\]

令\(Q,K\in\mathbb R^{n\times d_k}\)，\(V\in\mathbb R^{n\times d_v}\)，则注意力矩阵形状为\(n\times n\)，输出为\(n\times d_v\)。在自回归生成中，掩码\(M\)阻止当前token读取未来token。多头、残差、归一化与前馈层需要理解，但第一轮不必实现分布式训练。\href{https://arxiv.org/abs/1706.03762}{{[}P18{]}}
\href{https://courses.d2l.ai/zh-v2/}{{[}V03{]}}
\href{https://karpathy.ai/zero-to-hero.html}{{[}V10{]}}
\href{https://cs336.stanford.edu/spring2025/}{{[}V11{]}}

语言模型生成代码，实质上是在条件上下文下生成token序列。语法正确、格式遵从、算法可行和优化效果是逐级要求，不能从``next-token预测''直接推导出程序正确性。temperature改变采样分布，不是一个保证``创新程度''的物理旋钮；更高温度可能同时提高差异性和错误率。

\textbf{区分四件事。}
提示与上下文学习改变输入，不必改变模型权重；监督微调用示例程序更新权重；RL用奖励更新策略；外层进化筛选程序可以完全不更新模型权重。EoH初始复现不需要先掌握大模型微调。自然语言thought是候选的外显说明，不是对模型内部真实推理机制的证明。\href{https://arxiv.org/abs/2401.02051}{{[}P02{]}}
\href{https://arxiv.org/abs/2504.05108}{{[}P08{]}}
\href{https://www.youtube.com/watch?v=Ozos6M1JtIE}{{[}V07{]}}
\href{https://www.youtube.com/watch?v=Z6b5-77EfGk}{{[}V08{]}}
\href{https://www.youtube.com/watch?v=bJFtcwLSNxI}{{[}V09{]}}

\subsection{5.6
三个学习目标的数学区别}\label{ux4e09ux4e2aux5b66ux4e60ux76eeux6807ux7684ux6570ux5b66ux533aux522b}

\textbf{模仿学习。}
若示范数据为状态---动作对\((s,a^*)\)，离散动作策略可最小化\(\mathcal L_{\mathrm{IL}}(\theta)=-\mathbb E\log\pi_\theta(a^*\mid s)\)。用于路由时，状态可包含未服务任务和剩余容量，示范来自一个明确版本的专家算法；数据生成成本需要计入。训练时见过的专家状态与自己执行后到达的状态可能不同，低训练损失不等于低路线代价。\href{https://rail.eecs.berkeley.edu/deeprlcourse/}{{[}V14{]}}

\textbf{强化学习。}
令\(J(\theta)=\mathbb E_{\pi_\theta}[\sum_t\gamma^t r_t]\)，最简单的策略梯度估计使用\(\nabla_\theta\log\pi_\theta(a_t\mid s_t)\)乘以回报或优势估计。若奖励只在整条路线完成后给出，信用分配比监督一个下一站选择更困难。第一轮可用负路线长度作为终局奖励，但可行性仍由环境/验证器处理；不能仅靠一个任意罚项保证满足容量约束。\href{https://www.davidsilver.uk/teaching/}{{[}V13{]}}
\href{https://rail.eecs.berkeley.edu/deeprlcourse/}{{[}V14{]}}

\textbf{进化程序。}
EoH式搜索保留高分程序、再通过模型提出变化，通常没有上述对生成模型参数的梯度更新。课程学习则安排被呈现的训练分布，本身既不等于模仿损失，也不等于策略梯度。只有明确了这三者的对象，才有意义讨论怎样组合。\href{https://arxiv.org/abs/2401.02051}{{[}P02{]}}
\href{https://arxiv.org/abs/2504.05108}{{[}P08{]}}

\textbf{亲手练习。}
用三种手写策略作为动作，在少量训练实例上构造策略表现表；先实现总选一个策略的基线，再做监督选择。之后把它改成逐步选择算子的序列任务，写明奖励和状态变化。前者接近算法选择/上下文决策，后者才包含真正的时序信用分配；不要用同一张效果图掩盖任务变化。

\subsection{5.7
NAS与程序合成：不删除，但不夺走主线}\label{nasux4e0eux7a0bux5e8fux5408ux6210ux4e0dux5220ux9664ux4f46ux4e0dux593aux8d70ux4e3bux7ebf}

NAS需要明确搜索空间、搜索策略、性能估计和最终重训。DARTS通过连续松弛优化结构参数，是理解双层优化的好例子；它不是``LLM自动算法设计''的同义词。教学实验可限制为少量MLP结构与固定训练预算，比较随机搜索和一个简单搜索策略。先验证公平评估，再决定是否进入昂贵的图像NAS。\href{https://arxiv.org/abs/1806.09055}{{[}P14{]}}
\href{https://arxiv.org/abs/1808.05377}{{[}P15{]}}

程序合成从规格出发寻找程序。测试用例可以提供证据，但有限测试通常不构成对所有输入的证明。反例引导的思想是：提出候选、寻找违反规格的输入、用反例约束后续搜索。把它迁移到优化启发式时，必须区分``返回非法解''的正确性反例与``表现不佳''的性能反例；后者不代表程序逻辑错误。\href{https://people.csail.mit.edu/asolar/SynthesisCourse/index.htm}{{[}V18{]}}
\href{https://arxiv.org/abs/2108.07732}{{[}P17{]}}
\href{https://arxiv.org/abs/2403.03997}{{[}P24{]}}

\section{6
论文精读：把每篇读成一个可复现的实验}\label{ux8bbaux6587ux7cbeux8bfbux628aux6bcfux7bc7ux8bfbux6210ux4e00ux4e2aux53efux590dux73b0ux7684ux5b9eux9a8c}

\subsection{6.1
通用精读模板}\label{ux901aux7528ux7cbeux8bfbux6a21ux677f}

每篇论文先写一页卡片：研究对象是什么；被LLM改写的最小单位是什么；模型看见哪些数据；评分器返回什么；搜索/训练/推理预算各多少；最强的对照是谁；哪项消融能支持核心机制；有哪些无法重建的细节。最后补一个反例：``在哪种场景下，这个方法不应比简单基线更好？''

顺序应是摘要与图示定位问题，方法确定变量与接口，实验核对比较对象，再回到引言检查主张是否越界。不要用``这篇很强''``达到了SOTA''替代具体说明。每次阅读都留一个能运行或能核算的产物。

\subsection{6.2
EoH：第一篇深度复现对象}\label{eohux7b2cux4e00ux7bc7ux6df1ux5ea6ux590dux73b0ux5bf9ux8c61}

\textbf{核心理解。}
搜索对象是启发式，而不是一个问题实例的答案；候选携带自然语言说明和代码。探索与修改提示提出后代，真实评估决定保留。原论文包含在线装箱、TSP、流水车间调度；不要把后来仓库增加的CVRP等任务写进原论文实验范围。\href{https://arxiv.org/abs/2401.02051}{{[}P02{]}}

\textbf{必须重建的部分。}
用表格记录五种提示、父代抽样、初始种群、候选合法性处理、评估集、并列处理、停止条件与最大预算。原论文写每代每种策略产生一批候选时，不能用``代数相同''保证不同变体成本相同。你自己的现代模型复现需记录所有候选尝试和失败，而不仅是保留下来的程序。

\textbf{最有信息量的对照。}
同模型的独立采样、无thought的代码进化、固定初始候选的搜索、以及完整EoH。特别注意原论文中某些消融同时改变表示和提示集合；复现时需先忠实重建，再新增只改变一个因素的对照，不能把新实验反过来称为原论文设置。

\textbf{验收产物。}
一页方法说明、一份与论文/当前仓库差异表、一个可恢复运行的实验目录、全部候选的原始记录、冻结候选在独立测试集上的质量与失败统计。换了模型只能称``方法复现/现代模型迁移''，不能声称精确复现历史数字。

\subsection{6.3
FunSearch：评估闭环比语言流畅更重要}\label{funsearchux8bc4ux4f30ux95edux73afux6bd4ux8bedux8a00ux6d41ux7545ux66f4ux91cdux8981}

\textbf{核心理解。}
语言模型负责提出程序，评估器提供可计算信号，程序库与搜索过程持续保留有用候选。它研究cap
set和在线装箱，不是指南写的矩阵乘法系统。\href{https://deepmind.google/blog/funsearch-making-new-discoveries-in-mathematical-sciences-using-large-language-models/}{{[}P03{]}}
\href{https://github.com/google-deepmind/funsearch}{{[}C05{]}}

\textbf{阅读重点。}
找出人类固定了哪些程序骨架、输入规格和评价目标；检查候选如何进入程序库、怎样避免所有候选变得相似；说明``发现一个数学构造''和``得到对任意输入正确的算法''需要不同证据。

\textbf{复现边界。}
第一轮在同一个自建装箱评估器上实现FunSearch式对照，只能叫机制对照；要声称论文复现，还需对齐原任务分布、程序库机制、模型与预算。官方公开代码不自动包含你运行所需的全部服务和安全隔离。学习时先跑作者给出的已生成函数，再决定是否进行大规模生成。

\subsection{6.4
ReEvo：反思必须由行为证据约束}\label{reevoux53cdux601dux5fc5ux987bux7531ux884cux4e3aux8bc1ux636eux7ea6ux675f}

ReEvo将反思文本纳入进化过程，提供一个研究``反馈如何帮助后续生成''的清晰入口。读论文时区分短期/长期反馈的来源与作用，不把reflection自动等同于数值梯度。\href{https://proceedings.neurips.cc/paper_files/paper/2024/hash/4ced59d480e07d290b6f29fc8798f195-Abstract-Conference.html}{{[}P04{]}}
\href{https://github.com/ai4co/reevo}{{[}C06{]}}

可执行问题是：同样的候选对比，给分数、给简短失败摘要、给结构化反例，哪一种在相同token与候选评估预算下更有效？若复杂反思只是花更多token采样出更好的程序，不能把收益全部归于反思机制。还应检查反馈文本是否忠实于实际失败，不让LLM自行编造原因。

\subsection{6.5
LLM4AD：学习接口组织，不复制框架体量}\label{llm4adux5b66ux4e60ux63a5ux53e3ux7ec4ux7ec7ux4e0dux590dux5236ux6846ux67b6ux4f53ux91cf}

它把方法、任务与模型接口组织起来，是阅读现成系统的合适入口，但平台论文不是一个新的单一搜索算法。先用它找到任务评估器与方法调用点，再决定哪些设计可以借鉴。\href{https://arxiv.org/abs/2412.17287v2}{{[}P05{]}}
\href{https://github.com/Optima-CityU/LLM4AD}{{[}C03{]}}
\href{https://llm4ad-doc.readthedocs.io/en/latest/}{{[}C04{]}}

必须逐项核对README、\texttt{pyproject}或依赖文件、例子和安装版本。本次README要求Python不低于3.9且低于3.13；这不影响本册教学代码在另一Python版本通过测试，但不能据此宣布LLM4AD在该环境已验证。平台写``安全评估''也不免除你针对不可信生成代码的安全审查。

\subsection{6.6
AlphaEvolve与OpenEvolve：方法启发和原版复现分开}\label{alphaevolveux4e0eopenevolveux65b9ux6cd5ux542fux53d1ux548cux539fux7248ux590dux73b0ux5206ux5f00}

AlphaEvolve强调通过评估反馈持续改进代码，覆盖科学与工程问题；论文是白皮书，作者报告的具体成果应与它的任务、评估器和预算一同阅读。学习价值在于理解多种评估信号如何驱动代码改进，而不是复制一个产品名称。\href{https://arxiv.org/abs/2506.13131}{{[}P06{]}}

OpenEvolve是独立社区实现，能帮助你理解程序演化工程；它不是DeepMind原版源码，也不能仅凭名称相似就作为原版结果已复现的证据。第一轮不同时部署它、EoH和LLM4AD三个大系统，避免把接口调试当成方法研究。\href{https://github.com/algorithmicsuperintelligence/openevolve}{{[}C08{]}}

\subsection{6.7
必须纳入创新判断的近期研究}\label{ux5fc5ux987bux7eb3ux5165ux521bux65b0ux5224ux65adux7684ux8fd1ux671fux7814ux7a76}

EoH-S已把互补启发式集合本身作为设计对象；HeurAgenix涉及启发式生成和选择；FunTune把进化反馈用于RL微调生成器；ASRO预印本研究求解器与实例生成器共同演化；ES-AHD是本次查到的近期进化策略相关预印本。它们意味着``加集合''``加选择器''``加RL''``加课程学习''都不能直接作为未有人做过的创新主张。\href{https://ojs.aaai.org/index.php/AAAI/article/view/41038}{{[}P07{]}}
\href{https://arxiv.org/abs/2504.05108}{{[}P08{]}}
\href{https://arxiv.org/abs/2506.15196}{{[}P09{]}}
\href{https://arxiv.org/abs/2601.22896}{{[}P10{]}}
\href{https://arxiv.org/abs/2609.00023}{{[}P11{]}}

这里的核验深度并不相同：EoH是方法正文与示例层面的重点审读；近期条目主要用于确认先例和收窄问题，尚未独立复现实验。把预印本的作者主张、会议正式发表和本地实测结果分列，不混成一张确定无疑的排行榜。

\subsection{6.8
建议的最小必读集合与补充集合}\label{ux5efaux8baeux7684ux6700ux5c0fux5fc5ux8bfbux96c6ux5408ux4e0eux8865ux5145ux96c6ux5408}

主线先读EoH、FunSearch、ReEvo、LLM4AD，再读综述整理地图。讨论创新前加读EoH-S、FunTune及与你的机制直接相关的近期工作。Transformer、CodeGen、AlphaCode与NAS论文用于解释基础和边界；不需要把每篇历史模型论文都做全量实验。原指南的知识覆盖通过概念关卡完成，而不是靠把20篇论文都读到同一深度。

\section{7
视频、课程与社区：按问题观看，而不是按播放量收藏}\label{ux89c6ux9891ux8bfeux7a0bux4e0eux793eux533aux6309ux95eeux9898ux89c2ux770bux800cux4e0dux662fux6309ux64adux653eux91cfux6536ux85cf}

\subsection{7.1
资源核验方式与访问边界}\label{ux8d44ux6e90ux6838ux9a8cux65b9ux5f0fux4e0eux8bbfux95eeux8fb9ux754c}

本次检索覆盖作者和实验室主页、课程官网、arXiv与会议/期刊页面、GitHub官方与社区仓库、Hugging
Face课程、YouTube及B站的官方课程入口。它是有主题边界的资源筛选，不是``全网穷尽式系统综述''。保留的条目需要有明确教学作用；资源表记录一手来源、核验深度和访问限制。个人转述、营销式``最强教程''和无作者来源的搬运合集不承担技术证据。

资源编号只用于引用和查找，不是开发阶段代号。V类包括录像、课程、配套教材、讲义与媒体入口；\textbf{不是每个V编号都对应一门独立视频课}。同一报告的活动页与录像页分列，是为了提高可找回性，不重复计数。播放器可见、公开可下载、国内网络可访问和能否完成作业，是四个不同条件。本次没有逐帧观看全部视频，也没有实测你的校园网。

以下建议学时包含暂停推导和练习，不是视频官方时长；是可调整的学习预算。看不懂时先做最小例子，能够完成验收就跳过重复讲解，不用以``刷完整个课程''证明学习进度。

\Needspace{8\baselineskip}
\subsection{7.2
基础课程：直接服务本手册实验}\label{ux57faux7840ux8bfeux7a0bux76f4ux63a5ux670dux52a1ux672cux624bux518cux5b9eux9a8c}

\begin{longtable}[]{@{}
  >{\raggedright\arraybackslash}p{(\columnwidth - 4\tabcolsep) * \real{0.23}}
  >{\raggedright\arraybackslash}p{(\columnwidth - 4\tabcolsep) * \real{0.35}}
  >{\raggedright\arraybackslash}p{(\columnwidth - 4\tabcolsep) * \real{0.42}}@{}}
\toprule\noalign{}
\begin{minipage}[b]{\linewidth}\raggedright
学习问题
\end{minipage} & \begin{minipage}[b]{\linewidth}\raggedright
首选资源与观看范围
\end{minipage} & \begin{minipage}[b]{\linewidth}\raggedright
建议投入与验收
\end{minipage} \\
\midrule\noalign{}
\endhead
\bottomrule\noalign{}
\endlastfoot
怎样独立运行和排错 & Missing Semester
\href{https://missing.csail.mit.edu/}{{[}V02{]}}：Shell、Git、调试、性能分析
&
4---6小时；从干净目录重建环境，保存一次完整错误复现，不以截图代替日志 \\
为什么贪心不保证最优 & MIT 6.006
\href{https://ocw.mit.edu/courses/6-006-introduction-to-algorithms-spring-2020/}{{[}V01{]}}：图搜索、最短路、动态规划、复杂度
&
8---12小时选学；解释Dijkstra与TSP求解器为何解决不同问题，完成微型TSP穷举 \\
矩阵运算到底在做什么 & 3Blue1Brown
\href{https://www.3blue1brown.com/topics/linear-algebra}{{[}V15{]}}建立直觉，动手学深度学习
\href{https://courses.d2l.ai/zh-v2/}{{[}V03{]}}
\href{https://zh.d2l.ai/}{{[}V04{]}}做张量练习 &
4---6小时；不用自动广播``碰巧跑通''，手写每个矩阵维度 \\
随机实验如何比较 & STAT 110
\href{https://stat110.hsites.harvard.edu/}{{[}V16{]}}：随机变量、条件概率、期望与方差
& 6---8小时选学；区分一次搜索中的候选数与独立搜索重复数 \\
梯度从哪里来 & Karpathy
\href{https://karpathy.ai/zero-to-hero.html}{{[}V10{]}}中的micrograd，或中文D2L自动微分
& 6---10小时二选一；手算小计算图，再与自动微分核对 \\
约束和目标如何区分 & Boyd教材
\href{https://web.stanford.edu/~boyd/cvxbook/}{{[}V17{]}}：建模、约束、拉格朗日与对偶的入门部分
& 3---5小时；解释不可行解不能仅因路径短就获胜；无需先读完整本凸优化 \\
\end{longtable}

\textbf{观看顺序。}
Python能完成本册测试的人，从图算法与随机实验开始；不会建虚拟环境的人先做工具课。线性代数、概率和代码实践交错学习，不把三门完整课程堆在LLM实验之前。

\subsection{7.3
LLM与代码生成：中文主线、英文补充}\label{llmux4e0eux4ee3ux7801ux751fux6210ux4e2dux6587ux4e3bux7ebfux82f1ux6587ux8865ux5145}

\textbf{中文主线。}
用D2L的注意力章节解释公式与代码，再用李宏毅2025课程解释预训练、对齐、后训练和推理模型的差异。\href{https://courses.d2l.ai/zh-v2/}{{[}V03{]}}
\href{https://zh.d2l.ai/}{{[}V04{]}}
\href{https://speech.ee.ntu.edu.tw/~hylee/ml/2025-spring.php}{{[}V05{]}}
\href{https://www.youtube.com/watch?v=Xnil63UDW2o}{{[}V06{]}}
\href{https://www.youtube.com/watch?v=Ozos6M1JtIE}{{[}V07{]}}
\href{https://www.youtube.com/watch?v=Z6b5-77EfGk}{{[}V08{]}}
\href{https://www.youtube.com/watch?v=bJFtcwLSNxI}{{[}V09{]}}
第三讲偏内部机制剖析，不应代替从张量维度开始的注意力教学。第五至第七讲之后，写一张``模型参数有没有变化''的对照表：提示、检索、程序进化、监督微调、RL微调分别改变什么。

\textbf{实作入口。} Hugging
Face课程用于分词、模型加载、生成参数和数据流程；Karpathy课程用于看清小模型从数据到损失的过程。\href{https://karpathy.ai/zero-to-hero.html}{{[}V10{]}}
\href{https://huggingface.co/learn/llm-course/zh-CN/chapter1/1}{{[}V12{]}}
第一次生成实验只需要理解输入token、输出token、采样和解析，不以训练一个大模型作为前置条件。代码模型历史论文用来认识能力和评估边界，不强行部署指南中的所有旧模型。

\textbf{深入选修。}
CS336用于分词、训练细节和语言模型实验设计。\href{https://cs336.stanford.edu/spring2025/}{{[}V11{]}}
建议先用6---10小时选看相关课次；只有明确研究生成器训练的人，才追加大规模训练与系统优化作业。课程中使用的模型、工具版本和公开作业授权，应以所选学期页面为准。

\textbf{验收。}
给同一个任务建立可复用提示，保存原始输出及解析结果；解释为何温度变化会改变输出分布，但不会自动改变模型权重；能指出``返回了完整代码''与``代码通过独立评估''的差别。

\subsection{7.4
与自动算法设计最直接相关的四组材料}\label{ux4e0eux81eaux52a8ux7b97ux6cd5ux8bbeux8ba1ux6700ux76f4ux63a5ux76f8ux5173ux7684ux56dbux7ec4ux6750ux6599}

\textbf{EoH与LLM4AD。}
先读EoH方法正文，再看LLM4AD首页的TSP演示及作者提供的教程材料。\href{https://arxiv.org/abs/2401.02051}{{[}P02{]}}
\href{https://www.llm4ad.com/}{{[}V21{]}}
演示用于观察输入、生成、评估和结果的关系；CEC/EMO链接是教程幻灯片，不能因为名称叫tutorial就说有完整录像。第一轮用2---3小时串起论文图、README与任务接口。

\textbf{FunSearch作者报告。} Harvard
CMSA活动页能核对报告人、日期和题目，另有录像入口。\href{https://cmsa.fas.harvard.edu/event/mmlworkshop_1024/}{{[}V19{]}}
\href{https://www.youtube.com/watch?v=VJ_meFdGwE8}{{[}V20{]}}
观看时只记四件事：什么被固定、什么被生成、怎样评分、结果怎样验证。随后拿在线装箱做一次对应分析，避免停留在``AI发现数学''的故事层面。

\textbf{算法配置教程。} AutoML
2023的配置教程说明如何在给定参数空间内进行搜索；它正好能纠正``凡是让LLM调参数就是生成新算法''的混淆。\href{https://2023.automl.cc/program/tutorials/survey-of-methods/}{{[}V22{]}}
\href{https://www.tele-task.de/lecture/video/10161/}{{[}V23{]}}
选学3---4小时，建立一个不使用LLM的随机参数搜索，作为后续公平对照。

\textbf{语言模型与优化教程。} AutoML 2023的Language Modelling for
Optimization提供跨领域理解入口。\href{https://2023.automl.cc/program/tutorials/language-modelling-for-optimization/}{{[}V24{]}}
适合在完成第一次装箱闭环后观看；带着自己的输入/输出/预算表去对照，而不是第一次接触算法时直接进入宏观演讲。EoH-S官方页面还提供Video/Slides入口，用于后期讨论集合设计先例。\href{https://ojs.aaai.org/index.php/AAAI/article/view/41038}{{[}V25{]}}

\subsection{7.5
强化学习、模仿学习、课程学习何时进入}\label{ux5f3aux5316ux5b66ux4e60ux6a21ux4effux5b66ux4e60ux8bfeux7a0bux5b66ux4e60ux4f55ux65f6ux8fdbux5165}

David
Silver课程优先看MDP、价值函数和策略梯度；CS285优先看模仿学习及策略优化的相应课次。\href{https://www.davidsilver.uk/teaching/}{{[}V13{]}}
\href{https://rail.eecs.berkeley.edu/deeprlcourse/}{{[}V14{]}}
主线为理解预留8---12小时，不要求把深度强化学习全课修完。

最先要回答的是：\textbf{谁在学习，动作是什么，奖励什么时候给，训练后哪里发生变化？}
学习车辆下一站、学习选择局部搜索算子、学习生成程序，是三个不同问题。模仿学习可先模仿强启发式的选择；课程学习可改变实例难度或训练采样，但不能把测试集表现直接作为调课表的反馈。只有一个接口、一个训练对象和一个明确基线稳定后，才进入训练实验。

若研究只是生成一个离线启发式函数，RL并非必需。EoH原方法中的外部进化也不是RL微调；FunTune则提供了研究参数更新的相关先例。\href{https://arxiv.org/abs/2401.02051}{{[}P02{]}}
\href{https://arxiv.org/abs/2504.05108}{{[}P08{]}}
不为``融合更多方法''而增加无法单独验证的训练环节。

\subsection{7.6
社区使用方法：从提问到证据}\label{ux793eux533aux4f7fux7528ux65b9ux6cd5ux4eceux63d0ux95eeux5230ux8bc1ux636e}

对EoH/LLM4AD，优先查作者README、历史版本及Issues；对神经组合优化，查RL4CO的文档和仓库讨论；对程序进化，比较FunSearch原始代码与OpenEvolve社区实现。\href{https://github.com/FeiLiu36/EoH}{{[}C02{]}}
\href{https://github.com/Optima-CityU/LLM4AD}{{[}C03{]}}
\href{https://github.com/google-deepmind/funsearch}{{[}C05{]}}
\href{https://github.com/algorithmicsuperintelligence/openevolve}{{[}C08{]}}
\href{https://github.com/ai4co/rl4co}{{[}C17{]}}
社区答复是排错线索，不能替代论文实验或许可证文本。

一份有效提问包含：仓库与commit、Python及依赖、最小命令、完整异常、期望行为、已经排除的原因。不要把API密钥、付费账户信息、私有数据或完整机密日志贴进公开Issue。发现答案后在自己的实验记录中写下对应版本和变更，避免下次升级又遇到同一个问题。

每周最多预留1小时集中跟踪社区。新增资源必须说明它会改变哪一个实验决定；没有具体作用的收藏不进入主学习清单。资源失效时先回到作者机构/会议的上一级入口，再查归档，不使用来历不明的安装包替代。

\section{8
源码阅读与资源选择：先找到评估器，再找生成器}\label{ux6e90ux7801ux9605ux8bfbux4e0eux8d44ux6e90ux9009ux62e9ux5148ux627eux5230ux8bc4ux4f30ux5668ux518dux627eux751fux6210ux5668}

\subsection{8.1
仓库的职责与使用顺序}\label{ux4ed3ux5e93ux7684ux804cux8d23ux4e0eux4f7fux7528ux987aux5e8f}

\begin{longtable}[]{@{}
  >{\raggedright\arraybackslash}p{(\columnwidth - 4\tabcolsep) * \real{0.23}}
  >{\raggedright\arraybackslash}p{(\columnwidth - 4\tabcolsep) * \real{0.35}}
  >{\raggedright\arraybackslash}p{(\columnwidth - 4\tabcolsep) * \real{0.42}}@{}}
\toprule\noalign{}
\begin{minipage}[b]{\linewidth}\raggedright
资源
\end{minipage} & \begin{minipage}[b]{\linewidth}\raggedright
在本计划中的职责
\end{minipage} & \begin{minipage}[b]{\linewidth}\raggedright
第一遍看什么；暂时不做什么
\end{minipage} \\
\midrule\noalign{}
\endhead
\bottomrule\noalign{}
\endlastfoot
EoH \href{https://github.com/FeiLiu36/EoH}{{[}C02{]}} & 主方法复现入口 &
README版本提示、装箱/TSP例子、评估方法；不同时适配几十种任务 \\
LLM4AD \href{https://github.com/Optima-CityU/LLM4AD}{{[}C03{]}}
\href{https://llm4ad-doc.readthedocs.io/en/latest/}{{[}C04{]}} &
对照接口与统一任务组织 & 方法、任务、LLM接口；不先开发GUI和复杂服务 \\
FunSearch \href{https://github.com/google-deepmind/funsearch}{{[}C05{]}}
& 程序库与评估反馈的机制对照 &
程序骨架、评估和候选保留；不把社区复写视为原版结果 \\
ReEvo \href{https://github.com/ai4co/reevo}{{[}C06{]}} & 反思反馈对照 &
反馈来源、生成提示与演化过程；不只比较最终最好分数 \\
EoH-S \href{https://github.com/FeiLiu36/EoH-S}{{[}C07{]}} &
集合研究的必要先例 &
集合评价、互补性与测试流程；不直接宣称``集合''创新 \\
OpenEvolve / LLaMEA
\href{https://github.com/algorithmicsuperintelligence/openevolve}{{[}C08{]}}
\href{https://github.com/XAI-liacs/LLaMEA}{{[}C09{]}} &
扩展工程与优化器设计 & 选一个进行接口比较；不与主框架同时大改 \\
HeurAgenix \href{https://github.com/microsoft/HeurAgenix}{{[}C10{]}} &
生成和选择的组合先例 & 生成器、选择器及其训练/测试边界 \\
DEAP \href{https://github.com/DEAP/deap}{{[}C11{]}}
\href{https://deap.readthedocs.io/en/master/tutorials/basic/part2.html}{{[}C12{]}}
& 经典进化实现参照 & fitness方向、复制、变异后的失效与重评估 \\
PyVRP / HGS-CVRP \href{https://github.com/PyVRP/PyVRP}{{[}C13{]}}
\href{https://github.com/vidalt/HGS-CVRP}{{[}C14{]}} & 后续路由强基线 &
输入格式、成本精度、时间预算、版本；不是直接替代自己的问题定义 \\
ALNS \href{https://github.com/N-Wouda/ALNS}{{[}C16{]}} &
邻域搜索机制练习 &
destroy/repair、接受准则及算子选择；不预设为最终研究唯一骨架 \\
RL4CO \href{https://github.com/ai4co/rl4co}{{[}C17{]}} &
神经组合优化拓展 &
环境、策略、训练器、数据分布与评估；先跑小例子再考虑训练 \\
HumanEval / OR-Tools
\href{https://github.com/openai/human-eval}{{[}C18{]}}
\href{https://developers.google.com/optimization/routing/vrp}{{[}C15{]}}
& 正确性评估边界与路由参照 &
不可信代码安全警告、求解状态；通过代码测试不等于优化效果好 \\
\end{longtable}

当前PyVRP
README以迭代局部搜索描述求解器，不能把早期论文的HGS描述不加版本说明地套到当前代码。HGS-CVRP则是另一独立基线。选择谁取决于问题约束和预算，不取决于仓库名气。\href{https://github.com/PyVRP/PyVRP}{{[}C13{]}}
\href{https://github.com/vidalt/HGS-CVRP}{{[}C14{]}}

\subsection{8.2
原指南其余代码资源的去留}\label{ux539fux6307ux5357ux5176ux4f59ux4ee3ux7801ux8d44ux6e90ux7684ux53bbux7559}

BigCodeBench与APPS用于代码任务和功能正确性评估，不替代优化启发式benchmark；本次BigCodeBench页面已显示归档状态，应固定版本并评估依赖维护风险。\href{https://github.com/bigcode-project/bigcodebench}{{[}C19{]}}
\href{https://github.com/hendrycks/apps}{{[}C20{]}}
CodeGen、StarCoder和CodeLlama保留为历史模型与推理代码阅读入口，按模型卡和许可证决定是否使用，不强制全部部署。\href{https://github.com/salesforce/CodeGen}{{[}C21{]}}
\href{https://github.com/bigcode-project/starcoder}{{[}C22{]}}
\href{https://github.com/meta-llama/codellama}{{[}C23{]}}

NASBench的可核对入口为\path{google-research/nasbench}，它用于NAS基准学习；指南的泛化仓库名不能直接作为安装依据。\href{https://github.com/google-research/nasbench}{{[}C24{]}}
与AlphaCode相关的公开CodeContests仓库提供数据和相关材料，不等于完整AlphaCode训练系统的可复现发行版。\href{https://github.com/google-deepmind/code_contests}{{[}C25{]}}
这一组作为分支练习，主线仍是启发式生成、评估和搜索。

\subsection{8.3
真实可定位的入手点}\label{ux771fux5b9eux53efux5b9aux4f4dux7684ux5165ux624bux70b9}

\textbf{EoH。}
当前README说明：新API是\texttt{LLMConfig\ /\ EoH\ /\ BaseProblem}，旧\texttt{Paras\ /\ eoh.EVOL}对应旧版本。已核对的例子目录包括\path{examples/bp_online}、\path{examples/tsp_construct}，当前也有多函数、类形式与CVRP相关例子；这不等于原ICML论文已经研究了所有新增任务。\href{https://github.com/FeiLiu36/EoH}{{[}C02{]}}

\textbf{LLM4AD。}
已核对的README指向\path{example/online_bin_packing/online_bin_packing_tutorial.ipynb}、\path{example/tsp_construct/run_eoh.py}和\path{example/vrptw_construct/run_eoh.py}。其公开示例使用\texttt{OBPEvaluation}、\texttt{HttpsApi}、\texttt{EoH}及\texttt{EoHProfiler}。\href{https://github.com/Optima-CityU/LLM4AD}{{[}C03{]}}
本册不编造指南中不存在的\texttt{main.py\ -\/-problem\ obp}启动方式，也没有声称这些上游脚本在本环境运行通过。

\textbf{通用源码阅读顺序。}
从一个输入实例开始，追踪到评价结果：输入检查→人工基线→生成函数调用→独立验证→聚合指标→选择/保留→日志→下一次提示。只在必要时进入网络客户端、并行调度和可视化。每读完一段，用一句话说清它是否会影响科学结论；没有这个判断，逐行阅读也可能只是浏览。

\subsection{8.4
版本、环境与复现记录}\label{ux7248ux672cux73afux5883ux4e0eux590dux73b0ux8bb0ux5f55}

Ubuntu
22.04作为目标系统时，为教学代码、EoH/LLM4AD和神经模型训练分别建立虚拟环境。前两者先使用项目支持的Python版本；不要因系统Python或某个GPU包的要求，把全部仓库强行放进一个环境。指南并不要求一开始训练LLM。\href{https://github.com/FeiLiu36/EoH}{{[}C02{]}}
\href{https://github.com/Optima-CityU/LLM4AD}{{[}C03{]}}

每次克隆后先记录版本，再安装。下面是通用记录命令，\textbf{不是对所有仓库的统一安装脚本}：

\begin{verbatim}
git rev-parse HEAD
git status --short
python --version
python -m pip freeze > requirements-observed.txt
\end{verbatim}

复现实验应固定实际commit、数据校验值和配置，而不是长期跟踪\texttt{main}。README文件的blob
SHA只标识该文件内容，不能当成整个仓库的commit。锁文件、容器镜像和数据manifest各自解决不同问题，不能相互替代。升级后先跑相同的微型实例，再解释数值差异。

\subsection{8.5
算力与费用：先用可测量的预算说话}\label{ux7b97ux529bux4e0eux8d39ux7528ux5148ux7528ux53efux6d4bux91cfux7684ux9884ux7b97ux8bf4ux8bdd}

教学代码和经典小实例以CPU为主。调用远程LLM时，程序评估仍可在独立CPU环境执行；本地模型部署才需要进一步考虑显存和推理吞吐。8B参数以每参数2字节存权重，权重本身约16GB；这只是数量级计算，不包含KV缓存、激活、运行库及量化元数据，不能据此保证某张卡可运行完整任务。

不在本册固定易变的API价格、模型排行榜或并发额度。首次联网生成只开一个并发worker，设置硬预算上限；失败调用、重试、解析失败和无效候选都计入费用。先测20次调用的token、耗时及失败率，再决定扩大预算。这里的20是试运行的管理预算，不是统计充分性门槛。

手机适合阅读、标注、写问题清单及查看远程结果；完整编译、测试和安全执行放在已配置的电脑或云端工作目录。作者提供的Colab可以用于体验，但本次没有替你运行其笔记本，也不保证免费额度、网络连通性或第三方平台安全边界。不要为了得到一张演示图把密钥写进公开notebook。

\section{9
端到端最小系统：把可以信任的部分先建好}\label{ux7aefux5230ux7aefux6700ux5c0fux7cfbux7edfux628aux53efux4ee5ux4fe1ux4efbux7684ux90e8ux5206ux5148ux5efaux597d}

\subsection{9.1
系统的数据流与责任边界}\label{ux7cfbux7edfux7684ux6570ux636eux6d41ux4e0eux8d23ux4efbux8fb9ux754c}

\begin{longtable}[]{@{}
  >{\raggedright\arraybackslash}p{(\columnwidth - 4\tabcolsep) * \real{0.23}}
  >{\raggedright\arraybackslash}p{(\columnwidth - 4\tabcolsep) * \real{0.35}}
  >{\raggedright\arraybackslash}p{(\columnwidth - 4\tabcolsep) * \real{0.42}}@{}}
\toprule\noalign{}
\begin{minipage}[b]{\linewidth}\raggedright
环节
\end{minipage} & \begin{minipage}[b]{\linewidth}\raggedright
输入 → 输出
\end{minipage} & \begin{minipage}[b]{\linewidth}\raggedright
不应越过的边界
\end{minipage} \\
\midrule\noalign{}
\endhead
\bottomrule\noalign{}
\endlastfoot
问题与实例 & 任务规格、数据 → 验证后的实例 &
坐标、单位、约束和不可达关系必须明确 \\
人工算法与精确参照 & 实例 → 解及求解状态 &
微型精确结果与大规模最好已知结果分开 \\
任务包装器 & 实例、候选函数 → 完整解 &
可行性、终止和回退由可信代码负责 \\
候选生成 & 任务接口、训练反馈 → 候选源码 &
不向生成器暴露测试解、评分器写权限或密钥 \\
隔离执行 & 源码、受控输入 → 运行结果 &
独立受限环境；超时、崩溃和安全拒绝均留记录 \\
独立评估 & 完整解 → 正确性、可行性、质量、成本 &
不相信候选自报的目标值 \\
搜索与保留 & 训练评估 → 下一代/候选库 &
验证集只按事先规则用于少量选择 \\
测试与报告 & 冻结方法 → 测试结果与图表 &
查看测试结果后不得偷偷继续选择同一结论 \\
\end{longtable}

第一版只开放一个纯评分函数，例如\texttt{score(item,\ residuals)}或\texttt{priority(features)}。它返回分数，不自行决定文件访问、实验配置和结果写入。任务包装器用合法动作掩码、确定性的并列规则及回退策略生成完整解。接口简单不等于科研简化；它能使方法差异和失效来源更容易被识别。

\subsection{9.2
候选对象、运行记录与缓存}\label{ux5019ux9009ux5bf9ux8c61ux8fd0ux884cux8bb0ux5f55ux4e0eux7f13ux5b58}

一个候选至少记录：源码hash、父代hash、生成方法、提示hash、模型标识、调用参数、原始返回文本、解析状态及创建次序。一次评价至少记录：候选hash、实例hash、求解器版本、随机种子、预算、运行状态、独立目标值、可行性、时间和资源消耗。不要只保存最后一个``最佳.py''。

状态应能区分成功、不可行、超时、内存不足、语法错误、接口错误、崩溃、安全拒绝及API失败。API失败可能发生在候选产生之前，不能捏造成某个候选的差质量。聚合时同时报告过程可靠性与成功后的解质量。

缓存键至少覆盖代码、数据、配置、种子和评估器版本。完全相同的确定性评价可以缓存，但节省了多少计算要如实说明；随机评价不得因忽略seed被错误复用。代码文本重复与行为相同是不同层次的重复，后者需要预先规定的探针实例或行为特征。

\Needspace{18\baselineskip}
\subsection{9.3
搜索闭环的最小伪代码}\label{ux641cux7d22ux95edux73afux7684ux6700ux5c0fux4f2aux4ee3ux7801}

\begin{verbatim}
读取冻结的问题规格、训练集和总预算
用人工基线建立可运行的初始候选
当尚有预算时：
    从候选库选择父代或上下文
    生成一个候选；记录调用、重试与原始输出
    检查接口；在隔离环境运行微型安全/可行性测试
    若未通过，记录原因，不丢弃失败计数
    若通过，在训练实例上评价并计入预算
    按事先规定的保留规则更新候选库
按预定规则使用验证集选一个候选或一个固定集合
冻结源码、选择规则与配置
一次性运行保留测试集；导出全部状态与独立指标
\end{verbatim}

这段流程不是新的算法贡献，而是任何方法比较都需要的实验骨架。最初可用固定的手写候选代替网络生成器，验证整条数据流；这样的dry-run必须标为测试模式，不能计为LLM实验成果。

\subsection{9.4
生成代码的安全执行不是一个timeout参数}\label{ux751fux6210ux4ee3ux7801ux7684ux5b89ux5168ux6267ux884cux4e0dux662fux4e00ux4e2atimeoutux53c2ux6570}

AST检查、限制导入和子进程超时可以减少部分错误，却不是可靠的恶意代码安全边界。HumanEval官方仓库也明确警告，不可信代码不能未经隔离直接执行。\href{https://github.com/openai/human-eval}{{[}C18{]}}
学习中把候选视为不可信输入：主程序持有API密钥，评估worker不继承密钥；网络关闭；输入只读；可写临时目录有限；限制CPU、内存、进程数与运行时间；禁用特权和宿主机敏感挂载。

采用容器也不自动等于强安全；需要按风险选择额外隔离或虚拟机，并在运行未知代码前完成配置审查。此交付包\textbf{不实现生产级沙箱}，也没有包含运行任意LLM源码的\texttt{exec}入口。教学代码只运行随包附带、可人工审读的函数，避免把未验证的安全组件包装成科研能力。

\subsection{9.5
随包教学实验：已经运行了什么}\label{ux968fux5305ux6559ux5b66ux5b9eux9a8cux5df2ux7ecfux8fd0ux884cux4e86ux4ec0ux4e48}

\path{teaching/lab.py}与\texttt{test\_lab.py}提供最小纠错实验，使用本地NumPy与pytest，不需要API、外部数据或GPU。运行方法：

\begin{verbatim}
cd teaching
python -m pytest -q
python lab.py
\end{verbatim}

\begin{longtable}[]{@{}
  >{\raggedright\arraybackslash}p{(\columnwidth - 4\tabcolsep) * \real{0.23}}
  >{\raggedright\arraybackslash}p{(\columnwidth - 4\tabcolsep) * \real{0.35}}
  >{\raggedright\arraybackslash}p{(\columnwidth - 4\tabcolsep) * \real{0.42}}@{}}
\toprule\noalign{}
\begin{minipage}[b]{\linewidth}\raggedright
检查项
\end{minipage} & \begin{minipage}[b]{\linewidth}\raggedright
本次实测
\end{minipage} & \begin{minipage}[b]{\linewidth}\raggedright
能支持什么；不能支持什么
\end{minipage} \\
\midrule\noalign{}
\endhead
\bottomrule\noalign{}
\endlastfoot
单元测试 & 29项通过 & 教学函数通过这些测试；不是全面正确性证明 \\
原指南关键错误的最小重现 & \texttt{TypeError} &
新一代未评分便比较会出错；不是运行了原Word中的完整工程 \\
修正后的GA & 744次目标评估；历史最佳不升 &
评分顺序、预算计数与精英保留有效；不是某种优化器达到SOTA \\
正方形TSP & 最优4.0；交叉回路4.828427\ldots{} &
微型精确参照与独立成本核验正确 \\
装箱例子 & best-fit为2箱，next-fit为3箱 &
策略和实例影响质量；不能推导总体优劣 \\
\end{longtable}

本次教学环境为Python 3.13.5、NumPy 2.3.5、pytest
9.0.2；具体平台信息和日志随包保存。上游LLM4AD的版本要求不同，因此这些测试不能冒充其安装验证。GA数值只用于回归检查，不是跨平台位级一致性的承诺。

\subsection{9.6
从教学实验到论文复现的四次升级}\label{ux4eceux6559ux5b66ux5b9eux9a8cux5230ux8bbaux6587ux590dux73b0ux7684ux56dbux6b21ux5347ux7ea7}

先把教学装箱评估改成可载入固定实例集；再接入一个人工评分函数和一个固定候选库；随后接入一种真实LLM生成方法，记录所有调用；最后按论文设置建立基线、独立重复和测试分割。每次升级只增加一个影响结果的因素。任一步出错都能回退到已有可信结果，而不需要重建一个庞大平台。

\section{10
研究级实验协议：让改进经得住重新运行}\label{ux7814ux7a76ux7ea7ux5b9eux9a8cux534fux8baeux8ba9ux6539ux8fdbux7ecfux5f97ux4f4fux91cdux65b0ux8fd0ux884c}

\subsection{10.1
数据划分与泄漏控制}\label{ux6570ux636eux5212ux5206ux4e0eux6cc4ux6f0fux63a7ux5236}

至少区分设计训练集、方法选择验证集和最终测试集。训练集可用于反复进化；验证集只服务预先规定的超参数或候选选择；测试集在方法冻结后使用。若看了测试集又修算法，应建立新的独立测试证据，或如实把该集合降为开发数据。公开benchmark可能已经进入模型训练，不能声称仅靠随机切分就消除了预训练污染。

随机切分不是唯一方案。由同一地块、同一母实例或同一生成参数派生的样本应按组划分；防止旋转、尺度变化、点位扰动后的近重复实例跨集合。合成数据记录生成器、分布、seed与过滤条件，现实数据记录来源、许可、单位与预处理。删除失败或困难实例必须留理由，不能以``清洗''为名提高分数。

分布外测试至少分别考虑规模、几何/需求分布和约束变化，不把几种变化混成一个难以解释的OOD集合。规模外推不等于跨任务泛化；TSP到CVRP、CVRP到VRPTW也不是同一层次的问题。\href{https://arxiv.org/abs/2401.02051}{{[}P02{]}}
\href{https://ojs.aaai.org/index.php/AAAI/article/view/41038}{{[}P07{]}}
\href{https://github.com/ai4co/rl4co}{{[}C17{]}}

\subsection{10.2
必需基线和逐步增加的比较}\label{ux5fc5ux9700ux57faux7ebfux548cux9010ux6b65ux589eux52a0ux7684ux6bd4ux8f83}

\begin{longtable}[]{@{}
  >{\raggedright\arraybackslash}p{(\columnwidth - 4\tabcolsep) * \real{0.23}}
  >{\raggedright\arraybackslash}p{(\columnwidth - 4\tabcolsep) * \real{0.35}}
  >{\raggedright\arraybackslash}p{(\columnwidth - 4\tabcolsep) * \real{0.42}}@{}}
\toprule\noalign{}
\begin{minipage}[b]{\linewidth}\raggedright
要判断的主张
\end{minipage} & \begin{minipage}[b]{\linewidth}\raggedright
最少需要的对照
\end{minipage} & \begin{minipage}[b]{\linewidth}\raggedright
必须控制的因素
\end{minipage} \\
\midrule\noalign{}
\endhead
\bottomrule\noalign{}
\endlastfoot
LLM生成有效 & 人工规则、同预算独立LLM采样 &
模型、提示信息量、调用/评估成本 \\
进化优于重复采样 & 独立采样、固定候选库、完整进化 &
相同候选数与总时间；均报告失败 \\
自然语言反馈有效 & 仅分数、随机/打乱反馈、真实反馈 &
反馈token与生成预算，不能只控制迭代数 \\
设计比配置更有效 & 人工结构的随机/标准参数搜索 &
可调空间、试验数与开发成本 \\
集合/选择器有益 & 最佳单一算法、静态集合、随机选择 &
选择器额外成本及训练信息 \\
学习方法有益 & 手写策略、无训练对照、简单学习器 &
数据、训练和推理总成本 \\
农业约束有意义 & 同任务无该约束/不同强度的敏感性实验 &
不更换任务定义后直接比较总路长 \\
\end{longtable}

不是要求一次跑完所有基线，而是每提出一条主张补上对应证据。小问题先建立精确参照；大问题采用强公开求解器及明确的最好已知结果，不把未证明的参考值叫最优值。\href{https://github.com/PyVRP/PyVRP}{{[}C13{]}}
\href{https://github.com/vidalt/HGS-CVRP}{{[}C14{]}}
\href{https://developers.google.com/optimization/routing/vrp}{{[}C15{]}}

\subsection{10.3
预算：同时回答开发效率与使用效率}\label{ux9884ux7b97ux540cux65f6ux56deux7b54ux5f00ux53d1ux6548ux7387ux4e0eux4f7fux7528ux6548ux7387}

报告三条成本轴：模型生成的token与费用、候选评估次数与墙钟时间、冻结算法在新实例上的推理时间。相同迭代数不代表公平；失败重试和无效程序也是方法成本。并行硬件、线程数和终止策略一致或清晰标注。

若方法在部署前花费搜索成本\(C_{\mathrm{search}}\)，在\(M\)个实例上执行，每个实例开销为\(C_{\mathrm{run}}\)，则摊销成本是

\[
\bar C(M)=\frac{C_{\mathrm{search}}}{M}+C_{\mathrm{run}}.
\]

这不是给昂贵设计辩护，而是把``生成很贵、使用很快''的取舍写清。与参考法比较时，应代入两者各自的搜索/训练与使用成本，给出不同\(M\)下结论是否改变。一次性任务和重复作业任务可能得到不同推荐。

第一轮可规定20个候选用于接口试运行、100个候选用于方法排错；这些是管理上限而不是论文充分预算。正式实验预算根据初步方差、成本曲线和研究问题确定，不能先用测试成绩决定某个方法多跑多少轮。

\subsection{10.4
指标的定义与失败处理}\label{ux6307ux6807ux7684ux5b9aux4e49ux4e0eux5931ux8d25ux5904ux7406}

正成本且参考值大于零时，可报告相对差距

\[
g_i=\frac{f_i-f_i^{\mathrm{ref}}}{f_i^{\mathrm{ref}}}\times100\%.
\]

必须说明参考是精确最优、下界还是最好已知解；对零/负参考值使用预先定义的绝对差或其他归一化，不靠任意极小数掩盖问题。装箱等离散目标同时报告原始箱数；路径任务同时报告单位及成本定义。

可行性优先单独报告。只对成功实例计算均值会产生幸存者偏差；至少同时给出成功率、超时/不可行数量、共同成功实例的成对比较和对全部实例的明确处理。若使用惩罚值，必须预先说明依据，并给出敏感性分析，不能让某个随手选的罚常数决定结论。

候选级正确率与最终算法质量分开。候选失败率低可能来自保守搜索，不必然产生更好的解；最终最优分数高也可能来自极高试错成本。报告任一单指标都会漏掉方法的重要性质。

\subsection{10.5
随机性、独立单位与统计推断}\label{ux968fux673aux6027ux72ecux7acbux5355ux4f4dux4e0eux7edfux8ba1ux63a8ux65ad}

至少区分生成搜索seed、下层求解器seed、实例seed和训练seed。一个搜索过程产生1000个候选，仍不是1000次独立搜索。方法在同一测试实例上配对比较；多个运行共享训练数据的依赖也要说明。\href{https://jmlr.org/papers/v22/20-303.html}{{[}P22{]}}
\href{https://jmlr.org/papers/v7/demsar06a.html}{{[}P23{]}}

先进行小规模试验估计变异，再确定正式重复数。若预算只允许3---5次独立搜索，应报告分布和不确定性，不用``取最好一次''代替稳定性。作为可执行起点，本计划建议先做5次独立搜索的探索性比较；是否增至10次或更多由区间宽度与计算预算决定，不能预先宣称某个次数具有统计充分性。

报告效应量、成对差值和置信区间；bootstrap应按真正独立单位重采样，例如先按搜索运行、再按实例组分层，而不是把所有run-instance行当成独立样本。多数据集排名检验有适用条件，不能机械套用Demšar论文；多重比较和重复查看结果应提前考虑。实用改进门槛应来自应用容忍度或成本，而不是为了显著而事后指定。\href{https://jmlr.org/papers/v7/demsar06a.html}{{[}P23{]}}

\subsection{10.6
消融、停止与诚实的负结果}\label{ux6d88ux878dux505cux6b62ux4e0eux8bdaux5b9eux7684ux8d1fux7ed3ux679c}

第一篇报告只消融最关键的两个因素，例如反馈类型与多样性机制；其余保持固定。既跑``移除部件''，也验证简单替代方案能否达到相同效果。若同时换模型、提示、训练数据和搜索预算，结果不能归因于其中一个因素。

在正式运行前写下停止条件：预算耗尽、评估器错误、安全问题或预先约定的统计规则。安全/正确性问题触发立即停止；性能不够好不会触发偷偷换测试集。对失败方法写明适用范围、失效实例与成本，负结果也可以成为可靠的复现报告。

\subsection{10.7
应交付的图与表}\label{ux5e94ux4ea4ux4ed8ux7684ux56feux4e0eux8868}

推荐六类结果产物：质量---总成本曲线、各类失败比例、训练/验证/测试差距、规模及约束外推、关键机制消融、具体成功与失败案例。曲线展示多个独立运行的区间，不只画最好轨迹；对数轴、截断纵轴和归一化方式要解释。多目标任务可绘Pareto前沿，但只有目标存在实际取舍且单位/可行性明确时才采用，不能因图好看就增加目标。

图必须能从原始记录一键重建。自动生成的截图或任意模拟曲线只能标为示意，不进入结果区。本册没有伪造论文实验图；随包数值仅来自明确列出的教学运行。

\section{11
24周主线与前14天安排}\label{ux5468ux4e3bux7ebfux4e0eux524d14ux5929ux5b89ux6392}

\subsection{11.1
时间的使用原则}\label{ux65f6ux95f4ux7684ux4f7fux7528ux539fux5219}

主线按每周14小时安排，共336小时，是建议预算，不是``24周后一定能发表''的保证。每周用约4小时阅读/视频、7小时代码/实验、2小时结果检查、1小时写作与交流；临近复现实验可把观看时间转给运行分析。基础缺口允许前置0---8周；主线后另设8---12周研究扩展，不把扩展工作藏进同一份336小时计划。

不固定某天必须连续学习多久。每周至少保留两次不看参考实现的独立推导或编程，并保留一次从失败中解释原因的记录。章节、录像和代码共同围绕一个可验收问题，不要求每天同时学五门课。

\subsection{11.2
每周产物与过关条件}\label{ux6bcfux5468ux4ea7ux7269ux4e0eux8fc7ux5173ux6761ux4ef6}

\begin{longtable}[]{@{}
  >{\raggedright\arraybackslash}p{(\columnwidth - 4\tabcolsep) * \real{0.23}}
  >{\raggedright\arraybackslash}p{(\columnwidth - 4\tabcolsep) * \real{0.35}}
  >{\raggedright\arraybackslash}p{(\columnwidth - 4\tabcolsep) * \real{0.42}}@{}}
\toprule\noalign{}
\begin{minipage}[b]{\linewidth}\raggedright
周次
\end{minipage} & \begin{minipage}[b]{\linewidth}\raggedright
学习与实验任务
\end{minipage} & \begin{minipage}[b]{\linewidth}\raggedright
当周必须留下的证据
\end{minipage} \\
\midrule\noalign{}
\endhead
\bottomrule\noalign{}
\endlastfoot
第1周 & 阅读纠错章；运行教学包；补Shell/Git
\href{https://missing.csail.mit.edu/}{{[}V02{]}} &
环境记录、29项测试日志；口述原指南三项事实错误与一项代码错误 \\
第2周 & 贪心、图搜索与微型TSP
\href{https://ocw.mit.edu/courses/6-006-introduction-to-algorithms-spring-2020/}{{[}V01{]}}
& 独立实现/解释穷举参照；区分路径可达性与访问顺序 \\
第3周 & GA的选择、变异、重评估
\href{https://github.com/DEAP/deap}{{[}C11{]}}
\href{https://deap.readthedocs.io/en/master/tutorials/basic/part2.html}{{[}C12{]}}
& 单目标方向明确；预算计数；加入dim=1和零代数测试 \\
第4周 & 随机性、复杂度与公平比较
\href{https://stat110.hsites.harvard.edu/}{{[}V16{]}} &
对同一实例成对比较；说明为何候选数不是独立重复数 \\
第5周 & 注意力、token与自动微分
\href{https://courses.d2l.ai/zh-v2/}{{[}V03{]}}
\href{https://zh.d2l.ai/}{{[}V04{]}}
\href{https://karpathy.ai/zero-to-hero.html}{{[}V10{]}} &
手写形状推导；小计算图梯度核验；不依赖大型GPU \\
第6周 & 预训练、对齐、提示与采样
\href{https://speech.ee.ntu.edu.tw/~hylee/ml/2025-spring.php}{{[}V05{]}}
\href{https://www.youtube.com/watch?v=Xnil63UDW2o}{{[}V06{]}}
\href{https://www.youtube.com/watch?v=Ozos6M1JtIE}{{[}V07{]}}
\href{https://www.youtube.com/watch?v=Z6b5-77EfGk}{{[}V08{]}}
\href{https://www.youtube.com/watch?v=bJFtcwLSNxI}{{[}V09{]}}
\href{https://huggingface.co/learn/llm-course/zh-CN/chapter1/1}{{[}V12{]}}
& 一张参数是否更新表；一次有预算上限的生成记录，或明确的离线解析练习 \\
第7周 & 程序合成、规格、测试与NAS
\href{https://people.csail.mit.edu/asolar/SynthesisCourse/index.htm}{{[}V18{]}}
\href{https://arxiv.org/abs/1806.09055}{{[}P14{]}}
\href{https://arxiv.org/abs/1808.05377}{{[}P15{]}} &
小表达式合成练习；NAS与启发式设计异同比较；安全边界说明 \\
第8周 & 复盘三种教学任务，建立接口合同 &
装箱/TSP/CVRP微型问题规格；手写基线与独立验证，不急着训练 \\
第9周 & 精读EoH正文与当前README
\href{https://arxiv.org/abs/2401.02051}{{[}P02{]}}
\href{https://github.com/FeiLiu36/EoH}{{[}C02{]}} &
五个搜索提示算子的准确说明；原版/当前版差异表 \\
第10周 & 追踪一个装箱例子的完整调用链 &
一个实例从输入到fitness的追踪笔记；源码定位与版本记录 \\
第11周 & 接入受限生成与可靠性记录 &
小预算端到端运行；所有失败、重试和token保留；评估环境安全审查 \\
第12周 & 做EoH式首次方法复现 &
可恢复实验目录；与人工及独立采样对照；明确尚未对齐的论文设置 \\
第13周 & 学习FunSearch程序库机制
\href{https://deepmind.google/blog/funsearch-making-new-discoveries-in-mathematical-sciences-using-large-language-models/}{{[}P03{]}}
\href{https://github.com/google-deepmind/funsearch}{{[}C05{]}}
\href{https://cmsa.fas.harvard.edu/event/mmlworkshop_1024/}{{[}V19{]}} &
在同一包装器下的机制对照；与原版复现边界说明 \\
第14周 & 学习ReEvo反馈机制
\href{https://proceedings.neurips.cc/paper_files/paper/2024/hash/4ced59d480e07d290b6f29fc8798f195-Abstract-Conference.html}{{[}P04{]}}
\href{https://github.com/ai4co/reevo}{{[}C06{]}} &
分数/结构化反馈接口；打乱反馈对照设计 \\
第15周 & 读取LLM4AD接口
\href{https://arxiv.org/abs/2412.17287v2}{{[}P05{]}}
\href{https://github.com/Optima-CityU/LLM4AD}{{[}C03{]}} &
现有实现与平台的适配差异；不重新搭建一个完整GUI \\
第16周 & 冻结第一次比较协议，独立重复 &
数据分组划分、预算与运行seed清单；质量/成本/失败率表 \\
第17周 & 从TSP转向CVRP约束
\href{https://github.com/PyVRP/PyVRP}{{[}C13{]}}
\href{https://github.com/vidalt/HGS-CVRP}{{[}C14{]}}
\href{https://developers.google.com/optimization/routing/vrp}{{[}C15{]}}
& 小规模容量与回库核验；强基线输入对齐；不可达/不可行案例 \\
第18周 & 区分算法选择、配置、设计
\href{https://2023.automl.cc/program/tutorials/survey-of-methods/}{{[}V22{]}}
\href{https://2023.automl.cc/program/tutorials/language-modelling-for-optimization/}{{[}V24{]}}
& 同一人工结构的配置基线；说明自动生成到底新增什么自由度 \\
第19周 & RL与模仿学习入门
\href{https://www.davidsilver.uk/teaching/}{{[}V13{]}}
\href{https://rail.eecs.berkeley.edu/deeprlcourse/}{{[}V14{]}}
\href{https://github.com/ai4co/rl4co}{{[}C17{]}} &
为一个小问题定义状态/动作/奖励；简单模仿或表格学习练习 \\
第20周 & 课程与分布外泛化；查近期先例
\href{https://ojs.aaai.org/index.php/AAAI/article/view/41038}{{[}P07{]}}
\href{https://arxiv.org/abs/2504.05108}{{[}P08{]}}
\href{https://arxiv.org/abs/2506.15196}{{[}P09{]}}
\href{https://arxiv.org/abs/2601.22896}{{[}P10{]}}
\href{https://arxiv.org/abs/2609.00023}{{[}P11{]}} &
一个不使用测试集反馈的难度采样方案；创新重合检查表 \\
第21周 & 定义农业路径与任务问题，不先固定场景 &
单位、约束、地图、任务点与评估指标；导师可阅读的两页问题说明 \\
第22周 & 选择一条研究假设并做可行性试验 &
对照、预期证伪条件、算力预算；必要时放弃该假设而保留复现成果 \\
第23周 & 整理复现报告与可重建图表 &
方法、差异、结果、不确定性和失败案例；全部图表可追溯 \\
第24周 & 从干净环境复跑、口头答辩与修订 &
可运行打包产物；10分钟讲解；不能回答的问题进入明确补学清单 \\
\end{longtable}

第12、16、24周是建议的验收节点，不是冻结研究的行政门槛。若评估器不可信，第13周不进入新方法，而是继续修复；若基础已具备，可压缩课程周，把同等精力投入更多独立复现和失败分析。

\Needspace{8\baselineskip}
\subsection{11.3
前14天：把``开始学习''变成具体动作}\label{ux524d14ux5929ux628aux5f00ux59cbux5b66ux4e60ux53d8ux6210ux5177ux4f53ux52a8ux4f5c}

\begin{longtable}[]{@{}
  >{\raggedright\arraybackslash}p{(\columnwidth - 4\tabcolsep) * \real{0.23}}
  >{\raggedright\arraybackslash}p{(\columnwidth - 4\tabcolsep) * \real{0.35}}
  >{\raggedright\arraybackslash}p{(\columnwidth - 4\tabcolsep) * \real{0.42}}@{}}
\toprule\noalign{}
\begin{minipage}[b]{\linewidth}\raggedright
天
\end{minipage} & \begin{minipage}[b]{\linewidth}\raggedright
当天重点（约2小时）
\end{minipage} & \begin{minipage}[b]{\linewidth}\raggedright
完成标志
\end{minipage} \\
\midrule\noalign{}
\endhead
\bottomrule\noalign{}
\endlastfoot
1 & 阅读本册1---3章，对照原指南做更正 &
至少写清四角色、FunSearch和文献仓库三处问题 \\
2 & 运行教学测试并阅读日志 & 说明每类测试在保护什么，不能只说``29
passed'' \\
3 & 独立重现空适应度错误 & 写出先评分后选择的正确生命周期 \\
4 & 手算装箱序列，再改变到达次序 & 区分离线排序、在线决策和重新装箱 \\
5 & 手算正方形TSP并检查交叉回路 & 解释最优参照和独立计分器的用途 \\
6 & 看MIT图算法相关片段并画小图 &
区分最短路与TSP，不用最短路长度代替全任务成本 \\
7 & 不看源码讲解GA与装箱；写半页总结 &
讲出一个失败例子和一个尚不能解释的问题 \\
8 & 手写TSP置换验证器 & 重复节点、漏节点、越界输入均被发现 \\
9 & 给TSP做重标号和尺度测试 &
成本按预期变化，解的标签变化不造成伪改进 \\
10 & 画一个容量10、需求6/6/4的CVRP &
能判断哪些合并可行，理解回库不是错误 \\
11 & 阅读EoH摘要、方法图及任务说明 &
用自己的话说清``程序是个体，路线是解'' \\
12 & 看README并记录版本信息 & 识别新版API与旧教程；不盲装所有依赖 \\
13 & 写生成函数的输入输出合同 &
不允许候选修改评分器、数据划分或实验日志 \\
14 & 进行10分钟自测并修正前两周记录 &
能回答概念、代码、实验、公平性各一个问题 \\
\end{longtable}

第1---2周总量仍按28小时计；不要在这张日程之外再叠加一套14小时周任务。未接入真实模型前，离线解析和手写候选只能记作接口练习，不能写成完成LLM算法设计。

\subsection{11.4
基础补学与深入扩展}\label{ux57faux7840ux8865ux5b66ux4e0eux6df1ux5165ux6269ux5c55}

\textbf{需要前置补学时。}
用一组能力题诊断，而不是刷题数量：能否独立编写纯函数与测试；能否说明时间复杂度；能否解释期望和方差；能否读懂矩阵形状；能否从终端重跑一个程序。每个缺口补1---2周，最多按需要安排0---8周；已经会的项目直接免修。

\textbf{研究扩展8---12周。}
前2周补齐最相关先例及差异；随后2---3周实现唯一核心机制；再用2---4周做固定协议的重复与外推；最后2---3周复核、作图与写作。几个区间不是要求取所有上限，按总预算协调。研究没有出现可靠提升时，交付严谨复现、负结果和边界分析，不虚构创新。

\section{12
研究候选与农业规划迁移}\label{ux7814ux7a76ux5019ux9009ux4e0eux519cux4e1aux89c4ux5212ux8fc1ux79fb}

\subsection{12.1
不把技术名词拼接当成研究问题}\label{ux4e0dux628aux6280ux672fux540dux8bcdux62fcux63a5ux5f53ux6210ux7814ux7a76ux95eeux9898}

``LLM+RL+模仿学习+课程学习+VRP''只是技术组合，没有说明什么未知问题被解决。当前已有集合演化、启发式选择、生成器RL微调和求解器/实例共同演化的研究。\href{https://ojs.aaai.org/index.php/AAAI/article/view/41038}{{[}P07{]}}
\href{https://arxiv.org/abs/2504.05108}{{[}P08{]}}
\href{https://arxiv.org/abs/2506.15196}{{[}P09{]}}
\href{https://arxiv.org/abs/2601.22896}{{[}P10{]}}
合理起点是识别一个现有方法没有充分验证的机制，再用对照验证；不是先给系统起名、再寻找证明它有用的数据。

下面三个候选都是\textbf{待文献深入核查、待实验检验的设计假设}，不是声称已发现无人研究的方向。主线结束时只选一个，不同时进行三套研究。

\subsection{12.2
候选一：结构化反例反馈是否提高有效搜索效率}\label{ux5019ux9009ux4e00ux7ed3ux6784ux5316ux53cdux4f8bux53cdux9988ux662fux5426ux63d0ux9ad8ux6709ux6548ux641cux7d22ux6548ux7387}

\textbf{问题。}
与只有分数或自由文本反思相比，来源可核对的失败反例，能否在相同总预算下提高可行候选比例、降低找到有用启发式的成本？ReEvo和程序合成工作已经研究反馈与反例相关思想，因此创新不能只写``加入反思''。\href{https://proceedings.neurips.cc/paper_files/paper/2024/hash/4ced59d480e07d290b6f29fc8798f195-Abstract-Conference.html}{{[}P04{]}}
\href{https://arxiv.org/abs/2403.03997}{{[}P24{]}}

\textbf{最小实现。}
可信验证器把失败转成固定字段：违反的接口/约束、最小失败输入、候选实际输出、独立期望条件。失败缩减的计算也计入预算；生成器只看到训练实例中得到的反例。不可约束的自由文本建议留作对照。

\textbf{核心比较。}
同一模型、同一候选接口，比较只给分数、给长度匹配的通用说明、给随机/打乱反例、给真实结构化反例。分别报告正确性、可行性、质量与总成本；控制失败数量差异带来的反馈token变化。

\textbf{证伪条件。}
收益在token匹配后消失；只改善语法而不改善解质量；反例计算成本抵消搜索节省；换一种实例分布后没有收益。出现任一情况都需要缩小主张，而不是继续调图。

\subsection{12.3
候选二：带回退机制的启发式选择能否减少坏情况}\label{ux5019ux9009ux4e8cux5e26ux56deux9000ux673aux5236ux7684ux542fux53d1ux5f0fux9009ux62e9ux80fdux5426ux51cfux5c11ux574fux60c5ux51b5}

\textbf{问题。}
不是再证明一个集合优于单算法，而是研究低成本选择器何时应该放弃自信预测、回退到稳健基线，以减少分布变化时的严重退化。EoH-S与HeurAgenix是必须对照的先例。\href{https://ojs.aaai.org/index.php/AAAI/article/view/41038}{{[}P07{]}}
\href{https://arxiv.org/abs/2506.15196}{{[}P09{]}}

\textbf{最小实现。}
先固定一个小候选集合，提取求解前即可计算的实例特征，用简单树或线性模型预测相对表现；低置信度或超出训练范围时回退。特征开销、训练数据与回退阈值均由训练/验证确定，测试集不参加阈值选择。

\textbf{关键边界。}
未经形式化证明的经验回退不能叫``性能保证''；校准也不自动在任意OOD分布成立。至少对比最佳单算法、随机选择、始终回退、无回退选择器及不可部署的oracle上界。看最坏分位、尾部退化和平均成本，而不是只看平均路长。

\textbf{证伪条件。}
选择器费用高于改进；简单单算法已经足够；回退几乎从不触发或几乎一直触发；收益只来自在验证集反复搜索阈值。此时合理结论可能是该算法池无需复杂选择器。

\subsection{12.4
候选三：受控课程是否改善约束泛化}\label{ux5019ux9009ux4e09ux53d7ux63a7ux8bfeux7a0bux662fux5426ux6539ux5584ux7ea6ux675fux6cdbux5316}

\textbf{问题。}
在生成预算固定时，按规模、容量紧张度、障碍结构或时间窗紧张度安排训练实例，是否比固定混合采样更能产生跨分布有效的启发式？课程或共演化已有相关先例，必须与ASRO等区分；也不能把``从小到大''本身作为贡献。\href{https://arxiv.org/abs/2601.22896}{{[}P10{]}}

\textbf{最小实现。}
起初不训练实例生成器，只固定一个可解释的难度坐标和若干训练分布；比较均匀混合、固定递进和预先定义的自适应采样。难度应与实际可行性和求解成本相关，不能仅凭节点数命名``困难''。

\textbf{学习方法的合理位置。}
模仿学习可初始化算子选择器，RL可学习后续决策；但只有选定这一机制时才加入。先与不学习的课程比较，再研究学习课程是否值得；不同时更换启发式表示与生成模型。

\textbf{证伪条件。}
所谓泛化只发生在相似的生成器上；收益来自更多见过的实例而非课程；困难分布造成大面积无效候选；测试集参与调整课程。若这些问题无法解决，就保留均匀混合基线。

\subsection{12.5
农业任务如何进入，而不把主线带偏}\label{ux519cux4e1aux4efbux52a1ux5982ux4f55ux8fdbux5165ux800cux4e0dux628aux4e3bux7ebfux5e26ux504f}

农业迁移先写问题定义：机器人在哪种二维可通行空间中移动；任务是访问点、沿边服务、覆盖区域还是混合；任务需要多少作业资源；在哪里补给；存在什么时间或作业顺序约束；优化哪些具有实际意义的成本。只有这些条件成立，才能判断适合TSP、CVRP、VRPTW、弧路由、覆盖规划或其他模型，不能先选名词再解释场景。

\textbf{两个相连但不同的规划层。}
几何/图上的点到点规划提供可行转移与成本；任务规划决定访问哪些任务以及先后次序。若位置、朝向或入射方向影响下一段成本，单个静态距离矩阵可能不足，需要显式状态扩展或更合适的接口。矩阵元素对称、三角不等式成立以及所有点相互可达，都必须验证而不是默认。

\textbf{从仿真走向可信实例。}
第一批场景用可控二维地图，包含通道、禁行区、任务点和补给点，明确长度、时间和容量单位。再引入来源与许可可核对的数据，按原始场景分组划分训练/测试。若没有真实作业验证，只能声称二维模型下的方法效果，不能把地图演示称为实际农业部署。

\textbf{自动算法设计的位置。}
它可以改进候选动作评分、邻域排序、局部搜索调度或算法选择，但不能修改任务真值和评价规则来制造提升。几何层的错误不能通过更复杂LLM来掩盖。科研图首先展示地图与约束如何对应模型，再展示预算、公平对照及失败案例，不先追求``数字孪生''外观。

\section{13
从学习产物到可审查的论文}\label{ux4eceux5b66ux4e60ux4ea7ux7269ux5230ux53efux5ba1ux67e5ux7684ux8bbaux6587}

\subsection{13.1
论文结构与证据对应}\label{ux8bbaux6587ux7ed3ux6784ux4e0eux8bc1ux636eux5bf9ux5e94}

\begin{longtable}[]{@{}
  >{\raggedright\arraybackslash}p{(\columnwidth - 4\tabcolsep) * \real{0.23}}
  >{\raggedright\arraybackslash}p{(\columnwidth - 4\tabcolsep) * \real{0.35}}
  >{\raggedright\arraybackslash}p{(\columnwidth - 4\tabcolsep) * \real{0.42}}@{}}
\toprule\noalign{}
\begin{minipage}[b]{\linewidth}\raggedright
写作部分
\end{minipage} & \begin{minipage}[b]{\linewidth}\raggedright
要回答的问题
\end{minipage} & \begin{minipage}[b]{\linewidth}\raggedright
最低可追溯证据
\end{minipage} \\
\midrule\noalign{}
\endhead
\bottomrule\noalign{}
\endlastfoot
引言与相关工作 & 未解决的是哪一个具体问题 &
与最近先例逐机制比较；不能用``农业暂无应用''替代方法空白 \\
问题定义 & 输入、约束、输出、目标和假设是什么 &
单位、边界条件、可行性验证器和微型实例 \\
方法 & 新增了哪一步，为什么可能有效 &
伪代码、接口、预算计算；不把平台功能当成算法贡献 \\
实验设置 & 哪些条件固定，哪些允许调参 &
数据分割、模型、版本、种子、超参数和成本账单 \\
主要结果 & 效果大小、不确定性和成本是什么 &
全部实例结果、独立重复、强基线与失败统计 \\
消融与机制 & 收益是否真由该机制产生 &
单因素对照、简单替代方案、外推与案例 \\
局限和复现 & 结论在哪些情况下不成立 &
负结果、偏差来源、权限/安全边界、可运行材料 \\
\end{longtable}

不能先写``显著优于现有方法''再找实验支撑。开始写作时把主张分成已证明、实验支持、合理推测和未来工作四类；不同类型使用不同语气。顶刊要求是严谨性的目标，不是排版完成后的身份认证。\href{https://jmlr.org/papers/v22/20-303.html}{{[}P22{]}}

\subsection{13.2
投稿前必须能够回答的八个问题}\label{ux6295ux7a3fux524dux5fc5ux987bux80fdux591fux56deux7b54ux7684ux516bux4e2aux95eeux9898}

你是否能在不看代码时定义问题？评估器能否独立重算结果？原论文设置与改动是否透明？所有比较是否包含失败和完整成本？数据是否按真正的独立单位划分？新机制是否胜过简单替代？换分布后结论还剩什么？第三方是否能够用固定版本重建核心表格？

任何一项不能回答，都对应具体修复任务。工具数量、单元测试总数、代码行数、漂亮图和``多智能体''名称都不能代替这些答案。测试很多但验证错了目标，依然是不可信实验。

\subsection{13.3
每周汇报与答辩表达}\label{ux6bcfux5468ux6c47ux62a5ux4e0eux7b54ux8fa9ux8868ux8fbe}

每周汇报只需一页：本周问题、做过的实验、实际结果、发现的限制、下周唯一优先事项。明确``读过''``运行过''``复现过''``自己实现过''的区别。没有运行的数据就不填结果格；文献中的数字必须注明来源，不能混入个人实验表。

十分钟讲解按问题定义、方法机制、评估可信度、结果与局限组织。讲到老师指南中的错误时，可以说：``我沿用自动算法设计的主线，但核对原论文和仓库后修正了引文、角色分类与运行接口。''这是对材料的学术核查，不是对提供材料的人作价值判断。

\section{14
附录：概念自测与参考答案}\label{ux9644ux5f55ux6982ux5ff5ux81eaux6d4bux4e0eux53c2ux8003ux7b54ux6848}

\subsection{14.1
20个必须能独立回答的问题}\label{ux4e2aux5fc5ux987bux80fdux72ecux7acbux56deux7b54ux7684ux95eeux9898}

\begin{longtable}[]{@{}
  >{\raggedright\arraybackslash}p{(\columnwidth - 2\tabcolsep) * \real{0.5000}}
  >{\raggedright\arraybackslash}p{(\columnwidth - 2\tabcolsep) * \real{0.5000}}@{}}
\toprule\noalign{}
\begin{minipage}[b]{\linewidth}\raggedright
问题
\end{minipage} & \begin{minipage}[b]{\linewidth}\raggedright
合格答案的核心，不是背诵模板
\end{minipage} \\
\midrule\noalign{}
\endhead
\bottomrule\noalign{}
\endlastfoot
1. 一个TSP排列与一个启发式函数有何区别？ &
前者是某个实例的解；后者是可作用于多个实例的算法组成部分 \\
2. EoH中的个体是什么？ &
思想描述与可执行代码构成的启发式候选，不是必然为一条路线 \\
3. 综述四角色是什么？ &
optimizers、predictors、extractors、designers；不是指南的Programmer/Evolver四分法 \\
4. FunSearch与AlphaTensor有何区别？ &
前者是LLM程序搜索及外部评估；后者研究RL发现矩阵乘法算法，任务与方法不同 \\
5. EoH需要先微调一个LLM吗？ &
原方法不以微调为必需前提；外部进化不等于参数训练 \\
6. 为什么1000个候选不是1000次独立实验？ &
它们共享父代、反馈和搜索历史，存在依赖 \\
7. 为什么不能直接相信候选返回的cost？ &
候选可能算错、遗漏约束或操纵评分；可信验证器必须重算 \\
8. 最大化fitness能否优化最短路？ &
可以转换符号，但方向必须在所有选择/排序/报告中一致 \\
9. 训练最好为何不等于测试最好？ &
搜索反复适应训练样本，可能过拟合；选择还会引入偏差 \\
10. 相同代数为何不公平？ &
每代候选、重试、token、评估和并行资源可能不同 \\
11. first-fit、best-fit与next-fit有何区别？ &
选首个可行箱、最紧可行箱、只尝试当前箱；到达顺序会影响结果 \\
12. 微型TSP穷举有什么价值？ &
提供可核验的精确参照，发现成本/可行性错误；不解决大规模复杂性 \\
13. 距离矩阵何时可能不对称？ &
单行约束、方向状态或不同转移代价均可能导致方向相关；需建模核验 \\
14. 为什么HumanEval高分不足以证明启发式好？ &
编程正确性指标与优化目标、泛化、成本不是同一评价 \\
15. 一张GPU为何不代表能训练所有代码模型？ &
参数、优化器、梯度、激活、缓存和运行框架都有内存与计算成本 \\
16. 什么情况下应加入RL？ &
存在明确可学习的序列决策和可测奖励，而且简单基线不足以回答问题 \\
17. 课程学习会如何泄漏测试信息？ &
根据测试表现调整难度、生成器或采样规则会把测试变成开发反馈 \\
18. 超时能否当成缺失值直接删除？ &
不应；它是算法可靠性和预算结果的一部分，删除会偏向成功样本 \\
19. OpenEvolve是否是AlphaEvolve官方源码？ &
不是；它是社区实现，需按自身版本和任务评价 \\
20. 最终没有显著改进，还交付什么？ &
可信复现、成本/失败分析、适用边界、可重建数据与明确负结果 \\
\end{longtable}

\Needspace{28\baselineskip}
\subsection{14.2 概念索引}\label{ux6982ux5ff5ux7d22ux5f15}

实例instance：问题的一次具体输入；解solution：满足任务格式的输出；可行性feasibility：满足约束；目标objective：评价优劣的量；算法algorithm：从输入到输出的规则；启发式heuristic：利用结构快速求解的策略；配置configuration：给定结构的参数设置；选择selection：在已知方法间决策；设计design：搜索算法结构或代码。

种群population：并行维护的候选集合；父代parent：生成新候选的来源；变异mutation：产生局部/结构变化；交叉crossover：组合来源信息；精英保留elitism：保留已评价的优良候选；多样性diversity：避免候选过度同质；适应度fitness：搜索使用的评价，不必等同单个实例的目标值。

生成器generator：提出程序或数据的机制；评估器evaluator：执行并量化表现；验证器validator：独立检查格式、可行性或正确性；程序规格specification：程序应满足的要求；反例counterexample：违反规格的具体输入及行为；沙箱sandbox：限制不可信程序权限与资源的执行环境；接口合同contract：明确的输入输出和错误处理规定。

预训练pretraining：从大规模数据学习模型参数；微调fine-tuning：在后续数据/目标上更新参数；上下文学习in-context
learning：通过输入上下文改变行为；策略policy：状态到动作分布；回报return：累计奖励；模仿学习imitation
learning：从示范行为学习；课程curriculum：安排训练任务或采样难度；分布外OOD：与开发分布不同的评估条件。

独立重复independent
replication：使用独立随机过程重复完整方法；消融ablation：移除或替换机制以测试贡献；效应量effect
size：改进的大小；置信区间confidence
interval：量化估计不确定性的方法；预注册preregistration：事先写明假设和分析规则；工件artifact：代码、数据、配置和可重建结果；最好已知解BKS：当前已知参考，不自动是最优证明。

\clearpage
\section{15
参考文献与资源登记}\label{ux53c2ux8003ux6587ux732eux4e0eux8d44ux6e90ux767bux8bb0}

编号前缀P为论文与研究规范，V为视频、课程及配套材料，C为源码与工具，A为误引核查证据。共82条记录；包含相同报告的不同入口及7条误引证据，不等于82门课程或82个独立推荐。核验截止日期统一为2026年9月19日。原始指南为用户提供的非公开学习材料，本册不附其原文件。

\subsection{15.1
论文与研究规范}\label{ux8bbaux6587ux4e0eux7814ux7a76ux89c4ux8303}

\par\noindent\begin{minipage}{\linewidth}
\textbf{P01 · Liu et al.~A Systematic Survey on Large Language Models
for Algorithm Design. arXiv:\allowbreak{}2410.14716v5, 2026-01-02.}

入口：\url{https://arxiv.org/abs/2410.14716v5}

核验：一手页面已核验。按本次可读的 v5 引用；角色分类为 optimizers /
predictors / extractors / designers。
\end{minipage}\par\vspace{6pt}

\par\noindent\begin{minipage}{\linewidth}
\textbf{P02 · Liu et al.~Evolution of Heuristics: Towards Efficient
Automatic Algorithm Design Using Large Language Model. ICML 2024.}

入口：\url{https://arxiv.org/abs/2401.02051}

核验：一手页面已核验。方法正文及第3、4页截图核验；原论文任务为在线装箱、TSP、流水车间调度。
\end{minipage}\par\vspace{6pt}

\par\noindent\begin{minipage}{\linewidth}
\textbf{P03 · Romera-Paredes et al.~Mathematical discoveries from
program search with large language models. Nature 625, 468--475 (2024).
DOI:\allowbreak\nolinkurl{10.1038/s41586-023-06924-6.}}

入口：\url{https://deepmind.google/blog/funsearch-making-new-discoveries-in-mathematical-sciences-using-large-language-models/}

核验：作者机构、官方代码与出版社检索元数据已核验；出版社直接抓取受限。FunSearch；在线发表于2023年，卷期为2024年。
\end{minipage}\par\vspace{6pt}

\par\noindent\begin{minipage}{\linewidth}
\textbf{P04 · Ye et al.~ReEvo: Large Language Models as Hyper-Heuristics
with Reflective Evolution. NeurIPS 2024.}

入口：\url{https://proceedings.neurips.cc/paper_files/paper/2024/hash/4ced59d480e07d290b6f29fc8798f195-Abstract-Conference.html}

核验：一手页面已核验。
\end{minipage}\par\vspace{6pt}

\par\noindent\begin{minipage}{\linewidth}
\textbf{P05 · Liu et al.~LLM4AD: A Platform for Algorithm Design with
Large Language Model. arXiv:\allowbreak{}2412.17287v2, 2026-02-26.}

入口：\url{https://arxiv.org/abs/2412.17287v2}

核验：一手页面已核验。首版2024-12-23；平台论文不等于EoH原论文。
\end{minipage}\par\vspace{6pt}

\par\noindent\begin{minipage}{\linewidth}
\textbf{P06 · Novikov et al.~AlphaEvolve: A coding agent for scientific
and algorithmic discovery. arXiv:\allowbreak{}2506.13131, 2025.}

入口：\url{https://arxiv.org/abs/2506.13131}

核验：一手页面已核验。白皮书；不把作者结果写成本文复现结果。
\end{minipage}\par\vspace{6pt}

\par\noindent\begin{minipage}{\linewidth}
\textbf{P07 · Liu et al.~EoH-S: Evolution of Heuristic Set Using LLMs
for Automated Heuristic Design. AAAI 40(43):37090--37098, 2026.}

入口：\url{https://ojs.aaai.org/index.php/AAAI/article/view/41038}

核验：一手页面已核验。DOI:10.1609/aaai.v40i43.41038；预印本2508.03082。
\end{minipage}\par\vspace{6pt}

\par\noindent\begin{minipage}{\linewidth}
\textbf{P08 · Surina et al.~Algorithm Discovery With LLMs: Evolutionary
Search Meets Reinforcement Learning. arXiv:\allowbreak{}2504.05108, 2025.}

入口：\url{https://arxiv.org/abs/2504.05108}

核验：一手页面已核验。FunTune；搜索反馈用于RL微调生成器，不能称RL+进化全新。
\end{minipage}\par\vspace{6pt}

\par\noindent\begin{minipage}{\linewidth}
\textbf{P09 · Yang et al.~HeurAgenix: Leveraging LLMs for Solving
Complex Combinatorial Optimization Challenges. arXiv:\allowbreak{}2506.15196, 2025.}

入口：\url{https://arxiv.org/abs/2506.15196}

核验：一手页面已核验。本手册使用其启发式生成与选择机制进行先例核查；未运行上游实验。
\end{minipage}\par\vspace{6pt}

\par\noindent\begin{minipage}{\linewidth}
\textbf{P10 · Ke et al.~Game-Theoretic Co-Evolution for LLM-Based
Heuristic Discovery. arXiv:\allowbreak{}2601.22896, 2026.}

入口：\url{https://arxiv.org/abs/2601.22896}

核验：一手页面已核验。预印本；求解器/实例生成器共同演化，影响课程学习的创新判断。
\end{minipage}\par\vspace{6pt}

\par\noindent\begin{minipage}{\linewidth}
\textbf{P11 · Lai et al.~ES-AHD: An Evolution Strategy Framework for
Automatic Heuristic Design. arXiv:\allowbreak{}2609.00023, 2026.}

入口：\url{https://arxiv.org/abs/2609.00023}

核验：一手页面已核验。近期预印本，身份和摘要核验；未独立复现，不列为已证明的最强基线。
\end{minipage}\par\vspace{6pt}

\par\noindent\begin{minipage}{\linewidth}
\textbf{P12 · Li et al.~Competition-Level Code Generation with
AlphaCode. Science, 2022; arXiv:\allowbreak{}2203.07814.}

入口：\url{https://arxiv.org/abs/2203.07814}

核验：一手页面已核验。
\end{minipage}\par\vspace{6pt}

\par\noindent\begin{minipage}{\linewidth}
\textbf{P13 · Nijkamp et al.~CodeGen: An Open Large Language Model for
Code with Multi-Turn Program Synthesis. ICLR 2023; arXiv:\allowbreak{}2203.13474.}

入口：\url{https://arxiv.org/abs/2203.13474}

核验：一手页面已核验。
\end{minipage}\par\vspace{6pt}

\par\noindent\begin{minipage}{\linewidth}
\textbf{P14 · Liu, Simonyan, Yang. DARTS: Differentiable Architecture
Search. ICLR 2019; arXiv:\allowbreak{}1806.09055.}

入口：\url{https://arxiv.org/abs/1806.09055}

核验：一手页面已核验。
\end{minipage}\par\vspace{6pt}

\par\noindent\begin{minipage}{\linewidth}
\textbf{P15 · Elsken, Metzen, Hutter. Neural Architecture Search: A
Survey. JMLR 20(55), 2019; arXiv:\allowbreak{}1808.05377.}

入口：\url{https://arxiv.org/abs/1808.05377}

核验：一手页面已核验。
\end{minipage}\par\vspace{6pt}

\par\noindent\begin{minipage}{\linewidth}
\textbf{P16 · Jiang et al.~A Survey on Large Language Models for Code
Generation. arXiv:\allowbreak{}2406.00515, 2024.}

入口：\url{https://arxiv.org/abs/2406.00515}

核验：一手页面已核验。
\end{minipage}\par\vspace{6pt}

\par\noindent\begin{minipage}{\linewidth}
\textbf{P17 · Austin et al.~Program Synthesis with Large Language
Models. arXiv:\allowbreak{}2108.07732, 2021.}

入口：\url{https://arxiv.org/abs/2108.07732}

核验：一手页面已核验。研究论文，不冒充原指南所写的综述。
\end{minipage}\par\vspace{6pt}

\par\noindent\begin{minipage}{\linewidth}
\textbf{P18 · Vaswani et al.~Attention Is All You Need. NeurIPS 2017;
arXiv:\allowbreak{}1706.03762.}

入口：\url{https://arxiv.org/abs/1706.03762}

核验：一手页面已核验。
\end{minipage}\par\vspace{6pt}

\par\noindent\begin{minipage}{\linewidth}
\textbf{P19 · Rozière et al.~Code Llama: Open Foundation Models for
Code. arXiv:\allowbreak{}2308.12950, 2023.}

入口：\url{https://arxiv.org/abs/2308.12950}

核验：一手页面已核验。
\end{minipage}\par\vspace{6pt}

\par\noindent\begin{minipage}{\linewidth}
\textbf{P20 · Li et al.~StarCoder: may the source be with you!
arXiv:\allowbreak{}2305.06161, 2023.}

入口：\url{https://arxiv.org/abs/2305.06161}

核验：一手页面已核验。
\end{minipage}\par\vspace{6pt}

\par\noindent\begin{minipage}{\linewidth}
\textbf{P21 · Luo et al.~WizardCoder: Empowering Code Large Language
Models with Evol-Instruct. arXiv:\allowbreak{}2306.08568, 2023.}

入口：\url{https://arxiv.org/abs/2306.08568}

核验：一手页面已核验。
\end{minipage}\par\vspace{6pt}

\par\noindent\begin{minipage}{\linewidth}
\textbf{P22 · Pineau et al.~Improving Reproducibility in Machine
Learning Research. JMLR 22(164):1--20, 2021.}

入口：\url{https://jmlr.org/papers/v22/20-303.html}

核验：一手页面已核验。
\end{minipage}\par\vspace{6pt}

\par\noindent\begin{minipage}{\linewidth}
\textbf{P23 · Demšar. Statistical Comparisons of Classifiers over
Multiple Data Sets. JMLR 7:1--30, 2006.}

入口：\url{https://jmlr.org/papers/v7/demsar06a.html}

核验：一手页面已核验。统计思想需按实际独立实验单位使用，不能机械套用分类器检验。
\end{minipage}\par\vspace{6pt}

\par\noindent\begin{minipage}{\linewidth}
\textbf{P24 · Li, Parsert, Polgreen. Guiding Enumerative Program
Synthesis with Large Language Models. arXiv:\allowbreak{}2403.03997, 2024.}

入口：\url{https://arxiv.org/abs/2403.03997}

核验：一手页面已核验。用于理解LLM与形式化枚举搜索的互补，不是必须复现项目。
\end{minipage}\par\vspace{6pt}

\par\noindent\begin{minipage}{\linewidth}
\textbf{P25 · Fawzi et al.~Discovering faster matrix multiplication
algorithms with reinforcement learning. Nature 610, 47--53 (2022).
DOI:\allowbreak\nolinkurl{10.1038/s41586-022-05172-4.}}

入口：\url{https://deepmind.google/blog/discovering-novel-algorithms-with-alphatensor/}

核验：作者机构原始介绍已核验。AlphaTensor；用于与FunSearch区分。
\end{minipage}\par\vspace{6pt}

\subsection{15.2
视频、课程与讲义}\label{ux89c6ux9891ux8bfeux7a0bux4e0eux8bb2ux4e49}

\par\noindent\begin{minipage}{\linewidth}
\textbf{V01 · MIT 6.006 Introduction to Algorithms, Spring 2020.}

入口：\url{https://ocw.mit.edu/courses/6-006-introduction-to-algorithms-spring-2020/}

核验：一手页面已核验。课程主页含Video Lectures、讲义和习题；英文。
\end{minipage}\par\vspace{6pt}

\par\noindent\begin{minipage}{\linewidth}
\textbf{V02 · MIT The Missing Semester of Your CS Education, 2026 / 2020
archives.}

入口：\url{https://missing.csail.mit.edu/}

核验：一手页面已核验。Shell、Git、调试、性能分析；英文，站点列社区中文译本。
\end{minipage}\par\vspace{6pt}

\par\noindent\begin{minipage}{\linewidth}
\textbf{V03 · 李沐等：动手学深度学习在线课程，2021。}

入口：\url{https://courses.d2l.ai/zh-v2/}

核验：一手页面已核验。中文；作者课程页链接B站回放；单条B站播放器本次未验证。
\end{minipage}\par\vspace{6pt}

\par\noindent\begin{minipage}{\linewidth}
\textbf{V04 · 动手学深度学习中文版。}

入口：\url{https://zh.d2l.ai/}

核验：一手页面已核验。教材及可执行代码，不单独计为视频课程。
\end{minipage}\par\vspace{6pt}

\par\noindent\begin{minipage}{\linewidth}
\textbf{V05 · 李宏毅：Machine Learning 2025 Spring。}

入口：\url{https://speech.ee.ntu.edu.tw/~hylee/ml/2025-spring.php}

核验：一手页面已核验。中文；公开讲座链接可核验，部分校内作业平台不对外开放。
\end{minipage}\par\vspace{6pt}

\par\noindent\begin{minipage}{\linewidth}
\textbf{V06 · 李宏毅：第三讲，语言模型内部运作机制剖析。}

入口：\url{https://www.youtube.com/watch?v=Xnil63UDW2o}

核验：一手页面已核验。由V05官方课程页链接；不把这节替代完整Transformer数学教程。
\end{minipage}\par\vspace{6pt}

\par\noindent\begin{minipage}{\linewidth}
\textbf{V07 · 李宏毅：第五讲，预训练---对齐。}

入口：\url{https://www.youtube.com/watch?v=Ozos6M1JtIE}

核验：一手页面已核验。由V05链接；中文。
\end{minipage}\par\vspace{6pt}

\par\noindent\begin{minipage}{\linewidth}
\textbf{V08 · 李宏毅：第六讲，后训练与遗忘问题。}

入口：\url{https://www.youtube.com/watch?v=Z6b5-77EfGk}

核验：一手页面已核验。由V05链接；中文。
\end{minipage}\par\vspace{6pt}

\par\noindent\begin{minipage}{\linewidth}
\textbf{V09 · 李宏毅：第七讲，推理能力与DeepSeek-R1。}

入口：\url{https://www.youtube.com/watch?v=bJFtcwLSNxI}

核验：一手页面已核验。由V05链接；中文，课程内容为2025年。
\end{minipage}\par\vspace{6pt}

\par\noindent\begin{minipage}{\linewidth}
\textbf{V10 · Andrej Karpathy: Neural Networks: Zero to Hero.}

入口：\url{https://karpathy.ai/zero-to-hero.html}

核验：一手页面已核验。micrograd、makemore、GPT；英文，作者视频与代码。
\end{minipage}\par\vspace{6pt}

\par\noindent\begin{minipage}{\linewidth}
\textbf{V11 · Stanford CS336: Language Modeling from Scratch, Spring
2025.}

入口：\url{https://cs336.stanford.edu/spring2025/}

核验：一手页面已核验。选修分词、Transformer、实验；不要求完成大模型训练。
\end{minipage}\par\vspace{6pt}

\par\noindent\begin{minipage}{\linewidth}
\textbf{V12 · Hugging Face LLM Course，中文页面。}

入口：\url{https://huggingface.co/learn/llm-course/zh-CN/chapter1/1}

核验：一手页面已核验。中文讲义不代表所有视频为中文；接口实作参考。
\end{minipage}\par\vspace{6pt}

\par\noindent\begin{minipage}{\linewidth}
\textbf{V13 · David Silver: Reinforcement Learning course.}

入口：\url{https://www.davidsilver.uk/teaching/}

核验：一手页面已核验。官方教学页重定向到作者WordPress；MDP、策略梯度。
\end{minipage}\par\vspace{6pt}

\par\noindent\begin{minipage}{\linewidth}
\textbf{V14 · Berkeley Deep Reinforcement Learning course (CS285).}

入口：\url{https://rail.eecs.berkeley.edu/deeprlcourse/}

核验：一手页面已核验。从主页选择已结束学期；优先模仿学习与策略梯度。
\end{minipage}\par\vspace{6pt}

\par\noindent\begin{minipage}{\linewidth}
\textbf{V15 · 3Blue1Brown: Linear Algebra.}

入口：\url{https://www.3blue1brown.com/topics/linear-algebra}

核验：一手页面已核验。英文可视化直觉，不能替代矩阵形状推导。
\end{minipage}\par\vspace{6pt}

\par\noindent\begin{minipage}{\linewidth}
\textbf{V16 · Harvard STAT 110.}

入口：\url{https://stat110.hsites.harvard.edu/}

核验：一手页面已核验。概率论；选随机变量、条件概率、期望、极限定理。
\end{minipage}\par\vspace{6pt}

\par\noindent\begin{minipage}{\linewidth}
\textbf{V17 · Boyd \& Vandenberghe: Convex Optimization，作者资源页。}

入口：\url{https://web.stanford.edu/~boyd/cvxbook/}

核验：一手页面已核验。选约束与对偶；这是教材入口，不虚称完整免费视频。
\end{minipage}\par\vspace{6pt}

\par\noindent\begin{minipage}{\linewidth}
\textbf{V18 · Armando Solar-Lezama: Program Synthesis course.}

入口：\url{https://people.csail.mit.edu/asolar/SynthesisCourse/index.htm}

核验：一手页面已核验。MIT作者课程讲义；主要看规格、搜索、CEGIS。
\end{minipage}\par\vspace{6pt}

\par\noindent\begin{minipage}{\linewidth}
\textbf{V19 · Matej Balog: FunSearch，Harvard CMSA，2024-10-28。}

入口：\url{https://cmsa.fas.harvard.edu/event/mmlworkshop_1024/}

核验：一手页面已核验。主办方日程及YouTube播放列表已核验；不宣称逐帧观看。
\end{minipage}\par\vspace{6pt}

\par\noindent\begin{minipage}{\linewidth}
\textbf{V20 · Matej Balog \textbar{} FunSearch: Mathematical discoveries
from program search with LLMs.}

入口：\url{https://www.youtube.com/watch?v=VJ_meFdGwE8}

核验：视频索引与主办方事件交叉核验；直接播放器抓取受限。与V19为同一报告，不重复计为两门课。
\end{minipage}\par\vspace{6pt}

\par\noindent\begin{minipage}{\linewidth}
\textbf{V21 · LLM4AD: TSP video example, Colab, CEC/EMO tutorial links.}

入口：\url{https://www.llm4ad.com/}

核验：一手页面已核验。首页嵌入TSP演示；CEC/EMO链接为教程幻灯片，不冒充视频。
\end{minipage}\par\vspace{6pt}

\par\noindent\begin{minipage}{\linewidth}
\textbf{V22 · AutoML 2023: A Survey of Methods for Automated Algorithm
Configuration.}

入口：\url{https://2023.automl.cc/program/tutorials/survey-of-methods/}

核验：一手页面已核验。会议官方页面含Video；学习参数配置与算法设计边界。
\end{minipage}\par\vspace{6pt}

\par\noindent\begin{minipage}{\linewidth}
\textbf{V23 · AutoML 2023算法配置教程，tele-TASK录像。}

入口：\url{https://www.tele-task.de/lecture/video/10161/}

核验：主办方/录像索引核验；直接页面抓取受限。索引列英文及01:27:56；本手册不依赖播放时长排计划。
\end{minipage}\par\vspace{6pt}

\par\noindent\begin{minipage}{\linewidth}
\textbf{V24 · AutoML 2023: Language Modelling for Optimization.}

入口：\url{https://2023.automl.cc/program/tutorials/language-modelling-for-optimization/}

核验：一手页面已核验。Xingyou Song、Yutian Chen；概念桥接课。
\end{minipage}\par\vspace{6pt}

\par\noindent\begin{minipage}{\linewidth}
\textbf{V25 · AAAI 2026 EoH-S：Video/Slides。}

入口：\url{https://ojs.aaai.org/index.php/AAAI/article/view/41038}

核验：一手页面已核验。已核验页面存在Video/Slides按钮；未宣称媒体文件已下载。
\end{minipage}\par\vspace{6pt}

\subsection{15.3
源码、基准与文档}\label{ux6e90ux7801ux57faux51c6ux4e0eux6587ux6863}

\par\noindent\begin{minipage}{\linewidth}
\textbf{C01 · FeiLiu36/LLM4AlgorithmDesign（LLM4Opt）}

入口：\url{https://github.com/FeiLiu36/LLM4AlgorithmDesign}

核验：一手页面已核验。文献导航，不是运行框架。
\end{minipage}\par\vspace{6pt}

\par\noindent\begin{minipage}{\linewidth}
\textbf{C02 · FeiLiu36/EoH}

入口：\url{https://github.com/FeiLiu36/EoH}

核验：一手页面已核验。当前README：LLMConfig/EoH/BaseProblem；旧Paras/eoh.EVOL需v0.1；未运行上游。
\end{minipage}\par\vspace{6pt}

\par\noindent\begin{minipage}{\linewidth}
\textbf{C03 · Optima-CityU/LLM4AD}

入口：\url{https://github.com/Optima-CityU/LLM4AD}

核验：一手页面已核验。README blob
47e55f22d5058817a4e684339ee7b8273c74e013；不是commit；未运行上游。
\end{minipage}\par\vspace{6pt}

\par\noindent\begin{minipage}{\linewidth}
\textbf{C04 · LLM4AD documentation}

入口：\url{https://llm4ad-doc.readthedocs.io/en/latest/}

核验：官方README提供入口。安装/方法/任务/模型接口，版本需与仓库一致。
\end{minipage}\par\vspace{6pt}

\par\noindent\begin{minipage}{\linewidth}
\textbf{C05 · google-deepmind/funsearch}

入口：\url{https://github.com/google-deepmind/funsearch}

核验：一手页面已核验。算法结构学习；模型、沙箱、规模须单独核对。
\end{minipage}\par\vspace{6pt}

\par\noindent\begin{minipage}{\linewidth}
\textbf{C06 · ai4co/reevo}

入口：\url{https://github.com/ai4co/reevo}

核验：一手页面已核验。反思进化；未运行上游。
\end{minipage}\par\vspace{6pt}

\par\noindent\begin{minipage}{\linewidth}
\textbf{C07 · FeiLiu36/EoH-S}

入口：\url{https://github.com/FeiLiu36/EoH-S}

核验：一手页面已核验。互补集合设计；用于先例核查和进阶比较。
\end{minipage}\par\vspace{6pt}

\par\noindent\begin{minipage}{\linewidth}
\textbf{C08 ·
algorithmicsuperintelligence/openevolve（原codelion路径重定向）}

入口：\url{https://github.com/algorithmicsuperintelligence/openevolve}

核验：一手页面已核验。独立社区实现，不是DeepMind原版AlphaEvolve。
\end{minipage}\par\vspace{6pt}

\par\noindent\begin{minipage}{\linewidth}
\textbf{C09 · XAI-liacs/LLaMEA（原nikivanstein路径重定向）}

入口：\url{https://github.com/XAI-liacs/LLaMEA}

核验：一手页面已核验。优化器自动设计；先读README与评估配置。
\end{minipage}\par\vspace{6pt}

\par\noindent\begin{minipage}{\linewidth}
\textbf{C10 · microsoft/HeurAgenix}

入口：\url{https://github.com/microsoft/HeurAgenix}

核验：一手页面已核验。启发式生成与选择的既有工程。
\end{minipage}\par\vspace{6pt}

\par\noindent\begin{minipage}{\linewidth}
\textbf{C11 · DEAP}

入口：\url{https://github.com/DEAP/deap}

核验：一手页面已核验。进化操作及fitness生命周期练习。
\end{minipage}\par\vspace{6pt}

\par\noindent\begin{minipage}{\linewidth}
\textbf{C12 · DEAP: Operators and Algorithms}

入口：\url{https://deap.readthedocs.io/en/master/tutorials/basic/part2.html}

核验：一手页面已核验。clone和fitness失效必须理解；不能只照抄算子。
\end{minipage}\par\vspace{6pt}

\par\noindent\begin{minipage}{\linewidth}
\textbf{C13 · PyVRP/PyVRP}

入口：\url{https://github.com/PyVRP/PyVRP}

核验：一手页面已核验。当前README描述ILS；历史2024论文与最新主分支不是同一实现。
\end{minipage}\par\vspace{6pt}

\par\noindent\begin{minipage}{\linewidth}
\textbf{C14 · vidalt/HGS-CVRP}

入口：\url{https://github.com/vidalt/HGS-CVRP}

核验：一手页面已核验。README指向v1.0.0对应论文；读取round/seed/time设置。
\end{minipage}\par\vspace{6pt}

\par\noindent\begin{minipage}{\linewidth}
\textbf{C15 · Google OR-Tools: Vehicle Routing Problem}

入口：\url{https://developers.google.com/optimization/routing/vrp}

核验：一手页面已核验。模型与回调接口；启发式解不自动构成最优性证明。
\end{minipage}\par\vspace{6pt}

\par\noindent\begin{minipage}{\linewidth}
\textbf{C16 · N-Wouda/ALNS}

入口：\url{https://github.com/N-Wouda/ALNS}

核验：一手页面已核验。破坏/修复/选择/接受分离；不是预先指定最终研究框架。
\end{minipage}\par\vspace{6pt}

\par\noindent\begin{minipage}{\linewidth}
\textbf{C17 · ai4co/rl4co}

入口：\url{https://github.com/ai4co/rl4co}

核验：一手页面已核验。构造与改进策略；README含文档、Issues、Discussions、Slack。
\end{minipage}\par\vspace{6pt}

\par\noindent\begin{minipage}{\linewidth}
\textbf{C18 · openai/human-eval}

入口：\url{https://github.com/openai/human-eval}

核验：一手页面已核验。功能正确性；README明确警告执行不可信代码需要沙箱。
\end{minipage}\par\vspace{6pt}

\par\noindent\begin{minipage}{\linewidth}
\textbf{C19 · bigcode-project/bigcodebench}

入口：\url{https://github.com/bigcode-project/bigcodebench}

核验：一手页面已核验。页面显示2026-07-20归档；功能调用与复杂指令评估，未运行上游。
\end{minipage}\par\vspace{6pt}

\par\noindent\begin{minipage}{\linewidth}
\textbf{C20 · hendrycks/apps}

入口：\url{https://github.com/hendrycks/apps}

核验：GitHub连接器读取README已核验。APPS数据与训练/评价入口；网页直接抓取超时，但连接器成功。
\end{minipage}\par\vspace{6pt}

\par\noindent\begin{minipage}{\linewidth}
\textbf{C21 · salesforce/CodeGen}

入口：\url{https://github.com/salesforce/CodeGen}

核验：一手页面已核验。官方模型及程序合成入口；不作为2026最新推荐。
\end{minipage}\par\vspace{6pt}

\par\noindent\begin{minipage}{\linewidth}
\textbf{C22 · bigcode-project/starcoder}

入口：\url{https://github.com/bigcode-project/starcoder}

核验：一手页面已核验。官方微调与推理入口；运行与许可需按实际模型版本。
\end{minipage}\par\vspace{6pt}

\par\noindent\begin{minipage}{\linewidth}
\textbf{C23 · meta-llama/codellama}

入口：\url{https://github.com/meta-llama/codellama}

核验：一手页面已核验。推理代码入口；未下载权重，未验证当前设备兼容性。
\end{minipage}\par\vspace{6pt}

\par\noindent\begin{minipage}{\linewidth}
\textbf{C24 · google-research/nasbench}

入口：\url{https://github.com/google-research/nasbench}

核验：一手页面已核验。真实仓库名为nasbench；NAS数据与基准，非LLM启发式主线。
\end{minipage}\par\vspace{6pt}

\par\noindent\begin{minipage}{\linewidth}
\textbf{C25 · google-deepmind/code\_contests}

入口：\url{https://github.com/google-deepmind/code_contests}

核验：一手页面已核验。AlphaCode相关公开数据/材料，不等于完整原系统发行版。
\end{minipage}\par\vspace{6pt}

\subsection{15.4
原指南误引的核查证据}\label{ux539fux6307ux5357ux8befux5f15ux7684ux6838ux67e5ux8bc1ux636e}

\par\noindent\begin{minipage}{\linewidth}
\textbf{A01 · COWPEA (Candidates Optimally Weighted in Proportional
Election using Approval voting)}

入口：\url{https://arxiv.org/abs/2305.08857}

核验：一手页面已核验。原指南错配的真实标识符；仅用于查错，不是推荐阅读。
\end{minipage}\par\vspace{6pt}

\par\noindent\begin{minipage}{\linewidth}
\textbf{A02 · Probing Early Modification of Gravity with Planck, ACT and
SPT}

入口：\url{https://arxiv.org/abs/2308.12345}

核验：一手页面已核验。原指南错配的真实标识符；仅用于查错，不是推荐阅读。
\end{minipage}\par\vspace{6pt}

\par\noindent\begin{minipage}{\linewidth}
\textbf{A03 · Once-for-All: Train One Network and Specialize it for
Efficient Deployment}

入口：\url{https://arxiv.org/abs/1908.09791}

核验：一手页面已核验。原指南错配的真实标识符；仅用于查错，不是推荐阅读。
\end{minipage}\par\vspace{6pt}

\par\noindent\begin{minipage}{\linewidth}
\textbf{A04 · High-quality axions in a class of chiral U(1) gauge
theories}

入口：\url{https://arxiv.org/abs/2301.02345}

核验：一手页面已核验。原指南错配的真实标识符；仅用于查错，不是推荐阅读。
\end{minipage}\par\vspace{6pt}

\par\noindent\begin{minipage}{\linewidth}
\textbf{A05 · On Finite Analogs of Schmidt's Problem and Its Variants}

入口：\url{https://arxiv.org/abs/2205.00527}

核验：一手页面已核验。原指南错配的真实标识符；仅用于查错，不是推荐阅读。
\end{minipage}\par\vspace{6pt}

\par\noindent\begin{minipage}{\linewidth}
\textbf{A06 · The doubly librating Plutinos}

入口：\url{https://arxiv.org/abs/2501.12345}

核验：一手页面已核验。原指南错配的真实标识符；仅用于查错，不是推荐阅读。
\end{minipage}\par\vspace{6pt}

\par\noindent\begin{minipage}{\linewidth}
\textbf{A07 · 原指南LLM4AD链接：arXiv:\allowbreak{}2502.23456}

入口：\url{https://arxiv.org/abs/2502.23456}

核验：本次返回404。修正为P05，不据此断言一切同名资料不存在。
\end{minipage}\par\vspace{6pt}


\end{document}
